\documentclass[11pt]{article}
\usepackage[T1]{fontenc}
\usepackage{lmodern}
\usepackage{amsmath,amssymb,amsthm,mathtools}
\usepackage[margin=1in]{geometry}
\usepackage[colorlinks=true,linkcolor=blue,citecolor=blue,urlcolor=blue]{hyperref}
\usepackage{enumitem}
\allowdisplaybreaks
\setlength{\emergencystretch}{2em}
\numberwithin{equation}{section}
\newtheorem{theorem}{Theorem}[section]
\newtheorem{proposition}[theorem]{Proposition}
\newtheorem{lemma}[theorem]{Lemma}
\newtheorem{corollary}[theorem]{Corollary}
\theoremstyle{definition}
\newtheorem{definition}[theorem]{Definition}
\newtheorem{assumption}[theorem]{Imported input}
\theoremstyle{remark}
\newtheorem{remark}[theorem]{Remark}
\DeclareMathOperator{\Rea}{Re}
\DeclareMathOperator{\Ima}{Im}
\DeclareMathOperator{\Log}{Log}
\DeclareMathOperator{\erf}{erf}
\DeclareMathOperator{\erfi}{erfi}
\newcommand{\dd}{\,d}
\newcommand{\PV}{\operatorname{PV}}
\title{Signed Arithmetic Reductions on the Prime Torus\\
for Riemann Heat-Flow Collisions}
\author{Drafted for Edward Baker}
\date{10 October 2026}
\begin{document}
\maketitle
\begin{abstract}
We develop finite arithmetic reductions for a necessary fourth-jet test at
an ordinary real collision of the Riemann heat flow. The state uses the
prescribed quadratic logarithmic heat weights and one actual height
frequency; no independent prime phases are imposed. Complete centered
moments express the genuine finite candidates, raw-jet residuals, and
candidate-paid quadratic corrections. An extended correction makes every
ratio and product pair kernel vanish on the diagonal. A macroscopic pair
band of width comparable to the square root of the natural cutoff then has
a vanishing absolute payment. Large-common-factor difference coefficients
admit an Euler--Maclaurin transformation with a globally vanishing remainder.
For product terms we prove an oscillation-preserving finite
Poisson/incomplete-Fresnel formula with explicit endpoint transitions and
a globally paid error. An exact M\"obius decomposition keeps the primitive
ratio and product frontier coupled, with a uniform sublattice Poisson error.
The prescribed leading stationary reflection has small weight, center,
and carrier mismatches, but preserves the threshold quadratic and changes
the cutoff range. The required opposite signed bound for the retained
resonant and boundary terms remains open. The paper establishes reductions
and scoped obstructions, not a uniform collision exclusion or a proof of
the Riemann hypothesis.
\end{abstract}
\tableofcontents

\section{Problem, scope, and main reductions}
\label{sec:intro}
Project 09 investigates the arithmetic information visible in the complete
finite approximation to the Riemann heat flow. Its purpose is to replace
generic moment or independent-phase estimates by expressions that retain
the prescribed coefficients, the common height frequency, and both
ratio and product correlations.

The genuine heat function is normalized by
\[
H_t(z)=\int_0^\infty e^{tu^2}\Phi(u)\cos(zu)\dd u,
\qquad
\Phi(u)=\sum_{n\ge1}(2\pi^2n^4e^{9u}-3\pi n^2e^{5u})
                           e^{-\pi n^2e^{4u}},
\]
so that $\partial_tH_t=-H_t''$ and
$H_0(z)=\xi(1/2+iz/2)/8$. The classical threshold and the effective
approximation and zero count used here are recorded in
\cite{polymath}; this paper does not reprove those imported results.
At an all-real threshold an ordinary real double zero has a necessary
fourth-derivative inequality. After normalization, its paid arithmetic
counterpart is a quadratic in centered moments.
Our objective is an opposite inequality for the actual finite state.

The finite algebra below holds with the physical center, time, natural
integer cutoff, and selected coefficients frozen. Derivatives in
auxiliary summation variables are distinguished from genuine spatial
derivatives. Uniform asymptotic bounds mean all sufficiently small
positive times on the closed sector
$1\le\kappa\le3/2$; they do not supply estimates over complementary
parameters or the time-zero endpoint.

The principal analytical results have four components. First, a
candidate-paid degree-six extension cancels the diagonal separately in
each of the four pair kernels. Second, smoothing of difference
coefficients and a finite incomplete-Fresnel transformation of product
coefficients give explicit main terms with globally controlled errors.
Third, M\"obius inversion represents the small-common-factor frontier by
coupled sublattice quadratics without imposing candidate equations on the
sublattices. Fourth, exact physical reflection identities explain both
the useful near symmetry and the limitation of a leading stationary
reflection argument.

The proofs consolidate the project research notes \cite{projectnotes}
and the analytical heat continuations \cite{heatnotes}. Numerical controls
from the earlier scouts motivate some obstructions but are not used to
prove the analytical transformation theorems. No novelty in the literature
is asserted.

\section{The prescribed finite orbit and its raw jets}
\label{fo:state}

The finite state used here is the prescribed reflected approximation,
rather than a trigonometric polynomial with arbitrary coefficients.
For a complex spatial argument $z$, put $s=(1-iz)/2$ and, with
principal logarithms, define
\begin{align}
m_0(s)&=\Log s+\Log(s-1)-\frac{s}{2}\log\pi
 +\log\frac{\sqrt{2\pi}}{16}
 +\left(\frac{s}{2}-\frac12\right)\Log\frac{s}{2}-\frac{s}{2},
 \label{fo:m0}\\
\alpha(s)&=m_0'(s)=\frac1{2s}+\frac1{s-1}
 +\frac12\Log\frac{s}{2\pi},\qquad
m_t(s)=m_0(s)+\frac{t}{4}\alpha(s)^2.\nonumber
\end{align}
The relevant neighborhoods of the large positive real spatial axis
avoid the logarithmic cuts.  For a fixed positive integer $N$, let
\begin{align}
P_{t,N}(s)&=\sum_{n\le N}
 \exp\left\{\frac{t}{4}\log^2 n
       -\left(s+\frac{t}{2}\alpha(s)\right)\log n\right\},
 \label{fo:approximation}\\
G_{t,N}(z)&=e^{m_t(s)}P_{t,N}(s)
           +e^{m_t(1-s)}P_{t,N}(1-s),\nonumber\\
A_t(z)&=\exp\frac{m_t(s)+m_t(1-s)}2,\qquad
\theta_t(z)=\frac{m_t(s)-m_t(1-s)}{2i},\nonumber\\
F_{t,N}(z)&=G_{t,N}(z)/A_t(z),
 \qquad Q_t(z)=H_t(z)/A_t(z).\nonumber
\end{align}
Thus $A_t$ is nonvanishing, and on the real axis $A_t>0$ and
$F_{t,N},Q_t,\theta_t$ are real.  Consequently
\begin{equation}
H_t(x)=H_t'(x)=0\quad\Longleftrightarrow\quad
Q_t(x)=Q_t'(x)=0.
\label{fo:collision-equivalence}
\end{equation}

Throughout the analytical estimates we use the closed parameter sector
\begin{equation}
0<t\le\frac1{20},\qquad 1\le\kappa\le\frac32,\qquad
L=\frac{\kappa}{t},\qquad x=4\pi e^L,\qquad
N=\left\lfloor\sqrt{e^L+t/16}\right\rfloor.
\label{fo:sector}
\end{equation}
Every spatial derivative holds $t$ and this integer $N$ fixed.
The scale $L$ and the centering used below are the values at the chosen
center; neither is differentiated along \eqref{fo:sector}.
The complete fixed-cutoff disk estimate established for the parent
investigation will be used in the form
\begin{equation}
\left|Q_t(z)-F_{t,N}(z)\right|\le\eta_N
 \quad (|z-x|\le L^{-1}),\qquad
\eta_N\le5e^{-\mathfrak b/t},\qquad
\mathfrak b=\frac{\kappa(\kappa+4)}{16}.
\label{fo:disk-input}
\end{equation}
It includes conversion to $A_t$, reflection, and all possible
natural-cutoff changes in the disk.  This is an imported approximation
input \cite{polymath,heatnotes}, not a new smaller-cutoff theorem.
Cauchy's formula gives
\begin{equation}
\left|Q_t^{(j)}(x)-F_{t,N}^{(j)}(x)\right|
 \le j!L^j\eta_N\qquad(j\ge0).
\label{fo:cauchy}
\end{equation}

At the real center, write $\alpha=\alpha_r+i\alpha_i$ and
$\alpha'=U+iV$, where the prime on $\alpha$ means its complex
$s$ derivative.  Define
\begin{align}
\sigma&=\frac12+\frac{t\alpha_r}{2},&
w_n&=n^{-\sigma}e^{t\log^2 n/4},&
T&=\frac{x-t\alpha_i}{2},\label{fo:physical-data}\\
\phi_n&=\theta_t+T\log n,&
c&=\frac12\left(1+\frac{tU}{2}\right),&
d&=\frac{tV}{4},\nonumber\\
\Omega&=\frac12\operatorname{Re}\{\alpha(1+t\alpha'/2)\},&
\mu&=\frac{\Omega}{c},&
\epsilon&=\frac dc.\nonumber
\end{align}
The exact real-axis formulas are
\begin{align}
\alpha_r&=\frac L2+\frac14\log(1+x^{-2})-\frac1{1+x^2},&
\alpha_i&=\frac{3x}{1+x^2}-\frac12\arctan x,\label{fo:real-data}\\
U&=\frac{7x^2-5}{(1+x^2)^2},&
V&=\frac{x(x^2+5)}{(1+x^2)^2},\nonumber\\
\theta_t&=-\pi+\frac14\arctan x+\frac x4(1-L)
 -\frac x8\log(1+x^{-2})+\frac t2\alpha_r\alpha_i.\nonumber
\end{align}
In particular $c$ is uniformly bounded above and away from zero;
$0.49<c<0.51$ is a convenient reserve on \eqref{fo:sector}.
The physical sum is
\begin{equation}
q_n=w_ne^{i\phi_n},\qquad F_{t,N}(x)=2\operatorname{Re}\sum_{n\le N}q_n,
\qquad
M_j=\sum_{n\le N}q_n\rho_n^j=X_j+iY_j,\quad
\rho_n=\log n-\mu.
\label{fo:moments}
\end{equation}

If $\nu_p(n)$ is the exponent of $p$ in $n$, the same sum is a finite
prime-torus polynomial:
\begin{equation}
\sum_{n\le N}q_n
 =e^{i\theta_t}\sum_{n\le N}w_n
       \prod_{p\le N}\left(e^{iT\log p}\right)^{\nu_p(n)}.
\label{fo:torus}
\end{equation}
Unique factorization proves this identity without a density theorem.
All composites and all prescribed weights remain present.  Haar
orthogonality on the enlarged torus is an averaged statement; it is not
orthogonality at this one common height.  Moreover, for
$\mathcal D=\sum_{p\le N}\log p\,\partial_{\arg Z_p}$,
a monomial has $\mathcal D^2=-\log^2 n$.  The heat factor is generated
by $-\mathcal D^2/4$ on the finite monomial span.  Its positive-time
exponential on the full infinite span is unbounded.  The physical state
also moves $T$, its coefficients, and its carrier.

\subsection{Frozen centering and exact Bell rows}
\label{fo:raw-jets}
Direct differentiation of \eqref{fo:physical-data} gives
\begin{equation}
\frac{q_n'}{q_n}=-i\Omega+(-d+ic)\log n.
\label{fo:log-derivative}
\end{equation}
In a local differentiation calculation freeze $\mu$ and write
\begin{equation}
v_n=A+B\rho_n,\qquad A=-i\Omega+B\mu,\qquad B=-d+ic.
\label{fo:AB}
\end{equation}
At the center $A=a_{\rm c}:=-d\mu$ is real.  Its derivatives are
$A^{(h)}=-i\Omega^{(h)}+B^{(h)}\mu$ with $\mu$ frozen, and need not
be real.  Differentiating the center identity $A=-d\mu$ as though it
held on a moving centered chart would lose terms.

The raw multipliers are
\begin{align}
\mathcal B_0&=1,&\mathcal B_1&=v,\nonumber\\
\mathcal B_2&=v^2+v',&
\mathcal B_3&=v^3+3vv'+v'',\label{fo:Bell}\\
\mathcal B_4&=v^4+6v^2v'+3(v')^2+4vv''+v'''.\nonumber
\end{align}
Indeed differentiating $q_n\mathcal B_j$ proves
$\mathcal B_{j+1}=v_n\mathcal B_j+\mathcal B_j'$ and
$F_{t,N}^{(j)}/2=\operatorname{Re}\sum_n q_n\mathcal B_j(v_n)$.
For an explicit independent coefficient description let
$\mathcal B_j(v_n)=\sum_{k=0}^j p_{jk}\rho_n^k$.  Then
\begin{align}
p_{20}&=A^2+A',&p_{21}&=2AB+B',&p_{22}&=B^2,\nonumber\\
p_{30}&=A^3+3AA'+A'',&
p_{31}&=3A^2B+3(AB'+BA')+B'',\label{fo:Bell-rows}\\
p_{32}&=3AB^2+3BB',&p_{33}&=B^3,\nonumber
\end{align}
and
\begin{align}
p_{40}&=A^4+6A^2A'+3(A')^2+4AA''+A''',\nonumber\\
p_{41}&=4A^3B+6(A^2B'+2ABA')+6A'B'
                  +4(AB''+BA'')+B''',\nonumber\\
p_{42}&=6A^2B^2+6(2ABB'+B^2A')+3(B')^2+4BB'',
 \label{fo:fourth-row}\\
p_{43}&=4AB^3+6B^2B',\qquad p_{44}=B^4.\nonumber
\end{align}

The two observed lower coordinates are exactly
\begin{equation}
X_0=F_{t,N}/2,\qquad e_1=F_{t,N}'/2
   =a_{\rm c}X_0-dX_1-cY_1=a_{\rm c}X_0-cZ_1,
\qquad Z_1=Y_1+\epsilon X_1.
\label{fo:candidate-coordinates}
\end{equation}
At a genuine collision, \eqref{fo:cauchy} therefore implies
\begin{equation}
|X_0|\le a_0:=\eta_N/2,\quad |e_1|\le L\eta_N/2,
\quad |Z_1|\le a_1:=\frac{(L+|a_{\rm c}|)\eta_N}{2c}.
\label{fo:candidate-box}
\end{equation}
Neither $Y_0$ nor $X_1$ is independently zero.  For example, writing
$p_{j0}=u_0+iv_0$, $p_{j1}=u_1+iv_1$, exact elimination of $Y_1$
in \eqref{fo:Bell-rows}--\eqref{fo:fourth-row} yields
\begin{align}
F_{t,N}^{(j)}/2={}&-v_0Y_0+(u_1+dv_1/c)X_1
     +\operatorname{Re}\sum_{k=2}^j p_{jk}M_k\nonumber\\
&+(u_0-a_{\rm c}v_1/c)X_0+(v_1/c)e_1,
\qquad 2\le j\le4.
\label{fo:elimination}
\end{align}
The last line costs at most
$\eta_N|u_0-a_{\rm c}v_1/c|/2+L\eta_N|v_1|/(2c)$.
This identifies the actual hidden quadratures rather than replacing a
real value/derivative candidate by a complex zero of the torus sum.

\subsection{The exact finite heat residual}
\label{fo:heat-residual}
Let
\begin{equation}
b_n(s,t)=\exp\{m_0(s)-s\log n+t(\alpha(s)-\log n)^2/4\},
\quad h=\alpha-\log n,\quad B_s=1+t\alpha'/2.
\end{equation}
The logarithmic derivatives are $(\log b_n)_s=hB_s$ and
$(\log b_n)_{ss}=\alpha'B_s+(t/2)h\alpha''$.  Hence
\begin{equation}
\partial_t b_n-\frac14\partial_s^2b_n=R_n b_n,
\quad
R_n=\frac14\{h^2(1-B_s^2)-\alpha'B_s-(t/2)h\alpha''\}.
\label{fo:summand-residual}
\end{equation}
Using $\partial_x^2=-\partial_s^2/4$ on both branches gives
\begin{equation}
\partial_tG_{t,N}+G_{t,N}''=R_G,
\quad
R_G=\sum_{n\le N}\{R_n(s)b_n(s)+R_n(1-s)b_n(1-s)\}.
\label{fo:G-residual}
\end{equation}
With $\beta_A=A_t'/A_t$ and
$\varrho_A=A_t''/A_t+(\partial_tA_t)/A_t$, division by the same
normalizer gives
\begin{equation}
\partial_tF_{t,N}=-F_{t,N}''-2\beta_AF_{t,N}'-\varrho_AF_{t,N}
                   +R_G/A_t,\qquad
\partial_tQ_t=-Q_t''-2\beta_AQ_t'-\varrho_AQ_t.
\label{fo:normalized-residual}
\end{equation}
Already at $t=0$, $R_n=-\alpha'/4$.  The finite state is thus an
approximation to the scalar heat flow, not an invariant finite heat
manifold.  A spatial disk bound alone does not bound the time derivative
of its approximation error.

\section{A complete paid threshold reduction}
\label{fo:threshold-section}

\begin{lemma}[Uniform complete centered moments]
\label{fo:moment-lemma}
For every fixed nonnegative integer $j$, uniformly as $t\downarrow0$
in \eqref{fo:sector},
\begin{equation}
B_j:=\sum_{n\le N}w_n|\rho_n|^j\le C_jNw_N,
\quad \mu=\log N+O(N^{-1}),\quad
Nw_N\asymp e^{\mathfrak a/t},\quad
\mathfrak a=\frac{\kappa(4-\kappa)}{16}.
\label{fo:moment-bounds}
\end{equation}
The more general real-common-factor bound is
\begin{equation}
\sum_{a\le N/u}w(ua)(1+|\log(ua)-\mu|)^j
 \le C_j\frac Nu w_N\qquad(1\le u\le N),
\label{fo:real-index-moments}
\end{equation}
where $w(v)=v^{-\sigma}e^{t\log^2v/4}$.
\end{lemma}
\begin{proof}
The explicit data \eqref{fo:real-data} and the natural floor give
$\mu=\log N+O(N^{-1})$ and
$c_\sigma:=\sigma-(t/2)\log N=1/2+o(1)$.  They also give
$\log(Nw_N)=\mathfrak a/t+o(1)$, uniformly.  Set $A_u=N/u$ and
$y_a=\log(A_u/a)$.  The exact weight ratio is
\begin{equation}
w(ua)/w_N=\exp(c_\sigma y_a+t y_a^2/4).
\label{fo:weight-distance}
\end{equation}
For all sufficiently small time choose the uniform reserves
$7/16\le c_\sigma\le9/16$, $t\log N\le13/16$, and
$|\mu-\log N|\le1$.  Then
$c_\sigma+ty_a/4\le49/64$.  The bin $h\le y_a<h+1$ contains at
most $A_ue^{-h}$ integers.  In that bin the weight ratio is at most
$e^{49(h+1)/64}$ and the polynomial factor at most $(h+3)^j$.
Thus \eqref{fo:real-index-moments} holds with the admissible constant
\begin{equation}
C_j=e^{49/64}\sum_{h\ge0}e^{-15h/64}(h+3)^j<\infty.
\end{equation}
This counts the final low-index bin too.  Taking $u=1$ proves the
complete moment bound.  No endpoint term or macroscopic complement is
removed.
\end{proof}

For $j=2,3,4$, define the measured complete Bell residual
\begin{equation}
E_j=2\sum_{n\le N}w_n
       |\mathcal B_j(v_n)-(ic\rho_n)^j|.
\label{fo:Ej}
\end{equation}
For clarity put $b=ic\rho_n$ and $\xi=-d\log n$, so $v_n=b+\xi$
at the center.  Their exact residual polynomials are
\begin{align}
\mathcal B_2-b^2&=2b\xi+\xi^2+v',\nonumber\\
\mathcal B_3-b^3&=3b^2\xi+3b\xi^2+\xi^3+3(b+\xi)v'+v'',
 \label{fo:residual-polynomials}\\
\mathcal B_4-b^4&=4b^3\xi+6b^2\xi^2+4b\xi^3+\xi^4
 +6(b+\xi)^2v'+3(v')^2+4(b+\xi)v''+v'''.\nonumber
\end{align}
It follows, with no approximation to the coefficients, that
\begin{equation}
g_2=-2c^2X_2,\qquad g_3=2c^3Y_3,\qquad g_4=2c^4X_4,
\qquad |F_{t,N}^{(j)}-g_j|\le E_j.
\label{fo:leading-jets}
\end{equation}
The explicit data give $\xi=O(x^{-1})$ uniformly up to $N$, and
$v_n^{(h)}=O_h(x^{-h})$ for fixed $h\ge1$; these are fixed-time,
fixed-index derivatives.  Lemma~\ref{fo:moment-lemma} and
\eqref{fo:residual-polynomials} consequently show
\begin{equation}
E_j=O_j(Nw_N/x)=O_j(e^{\mathfrak a/t}/x),\qquad 2\le j\le4.
\label{fo:Ej-scale}
\end{equation}
The measured quantity used at a particular center remains
\eqref{fo:Ej}, rather than an unspecified asymptotic constant.

\subsection{The strengthened necessary sign}
\label{fo:count-threshold}
The following counting and canonical-product hypotheses are imported
from the genuine heat flow.  At an all-real time, opposite-root pairing
permits logarithmic differentiation of the canonical product.  Also an
absolute symbolic constant $C_{\rm count}\ge0$ may be fixed so that
\begin{equation}
|N_\tau(R)-p_\tau(R)|\le C_{\rm count}\log(2+R),
\quad 0<\tau\le\frac12,\quad R\ge4\pi,
\label{fo:count-input}
\end{equation}
where $N_\tau$ counts all zeros with $0<\operatorname{Re}\rho\le R$
including multiplicity, and
\begin{equation}
p_\tau(R)=\frac R{4\pi}\log\frac R{4\pi}-\frac R{4\pi}
             +\frac{11}{8}+\frac\tau{16}\log\frac R{4\pi}.
\end{equation}
The time $\tau$ in this input is distinct from the arithmetic
frequency $T$ \cite{polymath,heatnotes}.  No numerical value of
$C_{\rm count}$ is asserted.
Define
\begin{equation}
D_0=8\pi(4C_{\rm count}+1),\qquad
x_{\rm count}=\max\{(4\pi)^2,2D_0,3\},\qquad
\gamma=18(\log A_t)''+9/x^2,
\end{equation}
\begin{equation}
\gamma_{\rm count}=\gamma+\frac{9\log x}{2D_0^2},\qquad
\Gamma=\gamma_{\rm count}/c^2.
\label{fo:Gamma}
\end{equation}
All kernel algebra below is valid for any real frozen $\Gamma$;
the physical scale estimates use \eqref{fo:Gamma}.  In that case
$\gamma=O(x^{-2})$ and $\Gamma=\Theta(L)$ as $t\downarrow0$.

\begin{proposition}[Necessary ordinary-double threshold sign]
\label{fo:threshold-sign}
Assume $x\ge x_{\rm count}$ is an ordinary double zero at an all-real
time of the genuine even heat function.  Let $q_j=Q_t^{(j)}(x)$.
Then
\begin{equation}
2q_3^2-3q_2q_4-\gamma_{\rm count}q_2^2\ge0.
\label{fo:necessary-sign}
\end{equation}
\end{proposition}
\begin{proof}
Deflate the double root and write $h(z)=Q_t(z)/(z-x)^2$ locally.
For $S_0$ the inverse-square sum over the roots of $H_t$ remaining
after removing its two copies of $x$, logarithmic differentiation gives
\begin{equation}
S_0=\left(\frac{q_3}{3q_2}\right)^2
       -\frac{q_4}{6q_2}-(\log A_t)''.
\end{equation}
The mirror root $-x$ contributes $1/(2x^2)$.  On
$[x+D_0,x+2D_0]$ the derivative of $p_t$ is at least
$\log x/(8\pi)$ and each endpoint error in
\eqref{fo:count-input} is at most $2C_{\rm count}\log x$.
Therefore this disjoint interval contains at least
$(D_0/(8\pi)-4C_{\rm count})\log x=\log x$ roots counted with
multiplicity.  At the all-real time each supplies at least
$(2D_0)^{-2}$ to $S_0$.  Hence
\begin{equation}
S_0\ge\frac1{2x^2}+\frac{\log x}{4D_0^2}.
\end{equation}
Multiplication by $18q_2^2$ gives \eqref{fo:necessary-sign}.
The all-real hypothesis is essential in converting the counted roots
to a positive inverse-square sum.
\end{proof}

Let
\begin{equation}
\widehat\delta_j=j!L^j\eta_N+E_j,\qquad
\mathcal K=2Y_3^2+3X_2X_4-\Gamma X_2^2.
\label{fo:threshold-K}
\end{equation}
The complete measured quadratic perturbation payment is
\begin{align}
\widehat\Delta_{\rm count}={}&
4|g_3|\widehat\delta_3+2\widehat\delta_3^2
 +3\bigl(|g_2|\widehat\delta_4+|g_4|\widehat\delta_2
                         +\widehat\delta_2\widehat\delta_4\bigr)
 \nonumber\\
&+|\gamma_{\rm count}|
              (2|g_2|\widehat\delta_2+\widehat\delta_2^2).
\label{fo:Delta}
\end{align}
Indeed \eqref{fo:cauchy} and \eqref{fo:leading-jets} give
$|q_j-g_j|\le\widehat\delta_j$.  Expanding squares and products
proves the payment, and substitution gives
\begin{equation}
2g_3^2-3g_2g_4-\gamma_{\rm count}g_2^2=4c^6\mathcal K.
\label{fo:threshold-scaling}
\end{equation}
A sufficient exclusion of the collision in
Proposition~\ref{fo:threshold-sign} is consequently
\begin{equation}
\mathcal K+\widehat\Delta_{\rm count}/(4c^6)<0.
\label{fo:exclusion-target}
\end{equation}
This is the missing opposite arithmetic sign, not an established
inequality.  The strengthened coefficient must be included when
recomputing \eqref{fo:Delta}; the old payment with $\gamma$ cannot
be reused automatically.  The estimates above yield the admissible
uniform upper scale
\begin{equation}
\widehat\Delta_{\rm count}
 =O\left(L^4N^{-\kappa/4}+L^6N^{-(\kappa+4)/4}\right)=o(1).
\label{fo:Delta-scale}
\end{equation}
For example the products of a leading jet $O(e^{\mathfrak a/t})$
and a holomorphic jet error $O(L^je^{-\mathfrak b/t})$ have exponent
$\mathfrak a-\mathfrak b=-\kappa^2/8$, giving $N^{-\kappa/4}$.
Two holomorphic errors have exponent $-2\mathfrak b$, giving
$N^{-(\kappa+4)/4}$.  The new $O(L)$ normalizer factor contributes
at most $L^3$ and $L^5$ in these two cases and is absorbed by
\eqref{fo:Delta-scale}; Bell residual products are smaller.
The displayed scale is an upper budget, not a positive lower bound for
the measured payment.  A fourth-jet test can be vacuous at higher
multiplicity, so this ordinary-double criterion alone is not global
collision exclusion.

\section{Candidate-paid duals and the signed channels}
\label{fo:duals}

Set $e=(X_0,Z_1)^T$, $u=(X_2,Y_3,X_4)^T$, and
\begin{equation}
Q=\begin{pmatrix}-\Gamma&0&3/2\\0&2&0\\3/2&0&0\end{pmatrix},
\qquad \mathcal K=u^TQu.
\end{equation}
For a real symmetric $2\times2$ matrix $\Lambda$ and a real
$2\times3$ matrix $B$, chosen and frozen at the center, define
\begin{equation}
\mathcal K_{\Lambda,B}=\mathcal K+e^T\Lambda e+2e^TBu,
\qquad
\mathbb Q=\begin{pmatrix}\Lambda&B\\B^T&Q\end{pmatrix}.
\label{fo:dual-definition}
\end{equation}
The added expression vanishes on the exact finite candidate $e=0$.
On the paid box \eqref{fo:candidate-box}, its absolute value is at most
\begin{equation}
\Pi_{\Lambda,B}=
|\Lambda_{00}|a_0^2+2|\Lambda_{01}|a_0a_1+|\Lambda_{11}|a_1^2
 +2\sum_{j=0}^2(a_0|B_{0j}|+a_1|B_{1j}|)|u_j|.
\label{fo:box-dual-payment}
\end{equation}
Thus $\mathcal K_{\Lambda,B}\le U_{\Lambda,B}$ would be useful only
with the same parameters in the paid sufficient criterion
\begin{equation}
U_{\Lambda,B}+\Pi_{\Lambda,B}
                +\widehat\Delta_{\rm count}/(4c^6)<0.
\end{equation}
Large dual coefficients have no free margin.

Define the real feature vectors
\begin{equation}
r(\rho)=(1,\epsilon\rho,\rho^2,0,\rho^4)^T,
\qquad s(\rho)=(0,\rho,0,\rho^3,0)^T.
\end{equation}
For an ordered pair, let $R_{nm}=r_n^T\mathbb Qr_m$,
$S_{nm}=s_n^T\mathbb Qs_m$, and $C_{nm}=r_n^T\mathbb Qs_m$.
Product-to-sum proves exactly
\begin{align}
\mathcal K_{\Lambda,B}=\sum_{n,m\le N}w_nw_m\{&
D_c(n,m)\cos\Delta_{nm}+D_s(n,m)\sin\Delta_{nm}\nonumber\\
&+P_c(n,m)\cos\Sigma_{nm}+P_s(n,m)\sin\Sigma_{nm}\},
\label{fo:four-channels}\\
D_c&=(R+S)/2,\qquad P_c=(R-S)/2,\nonumber\\
D_s&=(C_{mn}-C_{nm})/2,\qquad
P_s=(C_{nm}+C_{mn})/2,\nonumber\\
\Delta_{nm}&=T\log(n/m),\qquad
\Sigma_{nm}=2\theta_t+T\log(nm).\nonumber
\end{align}
The difference-sine orientation follows from
$\sin(\phi_n-\phi_m)=\sin\phi_n\cos\phi_m-\cos\phi_n\sin\phi_m$.
Its coefficient and phase are both antisymmetric, so the ordered
sum does not cancel.  The prescribed weight factors as
\begin{equation}
w_nw_m=(nm)^{-\sigma}
  \exp\{t\log^2(nm)/8+t\log^2(n/m)/8\}.
\label{fo:pair-weight}
\end{equation}
Difference terms can be grouped by $n=ja,m=jb$ with $(a,b)=1$;
product terms can be grouped by $nm=k$ and its divisors.  These are
identities with one common $T$ and carrier, not independent Euler
factors.  For the physical block
$\mathcal B=\{\lfloor N/2\rfloor+1,\ldots,N\}$ and its complete
complement $\mathcal C$, every one of
$\mathcal B\times\mathcal B$, $\mathcal B\times\mathcal C$,
$\mathcal C\times\mathcal B$, $\mathcal C\times\mathcal C$
is retained by inserting its literal indicator inside the grouped sum.

For zero dual parameters both sine kernels vanish.  With
$h=\rho_n+\rho_m$ and $\delta=\rho_n-\rho_m$ the remaining kernels are
\begin{align}
D_c&=\frac{(h^2-\delta^2)^2(5h^2+\delta^2)}{128}
                  -\frac{\Gamma(h^2-\delta^2)^2}{32},\nonumber\\
P_c&=\frac{(h^2-\delta^2)^2(h^2+5\delta^2)}{128}
                  -\frac{\Gamma(h^2-\delta^2)^2}{32}.
\label{fo:base-kernels}
\end{align}
The leading coefficients are nonnegative while their cosine factors are
signed.  This does not supply the negative sign in
\eqref{fo:exclusion-target}.

\subsection{A sharp one-sided correlated payment}
\label{fo:curved-payment}
The disk error has a sharper first-jet body than the rectangle
\eqref{fo:candidate-box}.  At a collision apply Schwarz--Pick to
$g(\zeta)=(Q_t-F_{t,N})(x+\zeta/L)/\eta_N$ on the unit disk.
Real symmetry and $Q_t=Q_t'=0$ give the necessary condition
\begin{equation}
(2X_0/\eta_N)^2+2|a_{\rm c}X_0-cZ_1|/(L\eta_N)\le1.
\label{fo:curved-body}
\end{equation}
Indeed $|g'(0)|\le1-|g(0)|^2$.  Put
\begin{equation}
\mathcal C_0=\{(s,v):s^2+|v|\le1\},\qquad
\mathsf D=\begin{pmatrix}\eta_N/2&0\\
 a_{\rm c}\eta_N/(2c)&-L\eta_N/(2c)\end{pmatrix}.
\end{equation}
Then $e=\mathsf Dz$ for a $z\in\mathcal C_0$.  The sharp one-sided
payment on this necessary body is
\begin{equation}
\Pi^-_{\Lambda,B}(u)=
 -\min_{z\in\mathcal C_0}\{(\mathsf Dz)^T\Lambda\mathsf Dz
                            +2(\mathsf Dz)^TBu\}\ge0.
\label{fo:onesided}
\end{equation}
It bounds $\mathcal K-\mathcal K_{\Lambda,B}$ from above, and is
no larger than \eqref{fo:box-dual-payment}.  The sharpness here is
for the stated first-jet body, not a characterization of every
prescribed arithmetic state.

\begin{proposition}[Finite exact payment minimization]
\label{fo:cubic-payment}
Let $H=\mathsf D^T\Lambda\mathsf D=\left(\begin{smallmatrix}h&k\\k&r\end{smallmatrix}\right)$,
$g=\mathsf D^TBu=(l,m)^T$, and
$q(s,v)=hs^2+2ksv+rv^2+2ls+2mv$.
Its minimum on $\mathcal C_0$ is the minimum of its values at
$(\pm1,0)$, at every real derivative root in $(-1,1)$ of
\begin{equation}
P_\tau(s)=rs^4-2\tau ks^3+(h-2r-2\tau m)s^2
                +(2\tau k+2l)s+r+2\tau m,
\quad \tau=\pm1,
\label{fo:boundary-polynomial}
\end{equation}
and, if $\delta_H=hr-k^2\ne0$, at the feasible stationary point
\begin{equation}
s_*=(km-rl)/\delta_H,\quad v_*=(kl-hm)/\delta_H,
\quad q(s_*,v_*)=(-rl^2+2klm-hm^2)/\delta_H.
\label{fo:payment-stationary}
\end{equation}
An identically constant boundary is already covered by its endpoints.
No additional interior value is necessary when $H$ is singular.
\end{proposition}
\begin{proof}
The two boundary arcs are $v=\tau(1-s^2)$, and substitution gives
\eqref{fo:boundary-polynomial}; their critical equations are cubic.
An interior minimum satisfies $Hz+g=0$, which gives
\eqref{fo:payment-stationary} in the nonsingular case.  If $H$ is
singular and a stationary point exists, the objective is constant on
a line through that point in a nonzero null direction of $H$.
The line reaches the boundary of the compact body with the same value.
This includes stationary ridges and the identically zero quadratic.
\end{proof}
For a complete independent moment box, the payment is exactly the
maximum of \eqref{fo:onesided} over its at most eight common corners:
for fixed $z$ the objective is affine in $u$, so the two maxima can be
interchanged.  Likewise an independent interval box for
$(h,k,r,l,m)$ can be paid at its at most 32 common corners.
These are exact finite optimization facts; actual correlations lost
when forming a box can improve the payment further.

\subsection{What norm relaxation and Gaussian averaging retain}
\label{fo:relaxations}
The threshold matrix $Q$ has two positive and one negative eigenvalue:
its $Y_3$ diagonal is $2$ and its $X_2,X_4$ block has determinant
$-9/4$.  This limits a precise candidate-conditioned norm relaxation.
To state it, put $W=\sum_nw_n$, $p_n=w_n/W$, and
$R_*^2=\sum_np_n\rho_n^2>0$.  In the real Hilbert space with weights
$p_n$, use the two lower shapes
\begin{equation}
a_n=\cos\phi_n,\quad
b_n=(\rho_n/R_*)(\sin\phi_n+\epsilon\cos\phi_n)
\end{equation}
and the three upper shapes
$f_2=(\rho/R_*)^2\cos\phi$,
$f_3=(\rho/R_*)^3\sin\phi$,
$f_4=(\rho/R_*)^4\cos\phi$.  Let $G$ be the lower Gram matrix,
$C$ the lower--upper Gram block, $F$ the upper Gram matrix, and
$H=F-C^TG^{-1}C$ when $G$ is invertible.  If
\begin{equation}
z=\left(X_0/W,Z_1/(WR_*)\right)^T,\quad
m=\left(X_2/(WR_*^2),Y_3/(WR_*^3),X_4/(WR_*^4)\right)^T,
\quad v=C^TG^{-1}z,
\end{equation}
projection and Cauchy--Schwarz give
\begin{equation}
(m-v)(m-v)^T\preceq(1-z^TG^{-1}z)H.
\label{fo:Gram-relaxation}
\end{equation}
To verify it against any vector of upper coefficients, project the
constant function and the corresponding upper shape off the lower
span.  Their squared residual norms are respectively
$1-z^TG^{-1}z$ and the appropriate quadratic form in $H$.
At an exact candidate $z=0$, the sharp norm-ball maximum is
\begin{equation}
\sup_{mm^T\preceq H}m^TQ_*m
 =\max\{0,\lambda_{\max}(H^{1/2}Q_*H^{1/2})\},\qquad
Q_*=\begin{pmatrix}-\Gamma/R_*^2&0&3/2\\0&2&0\\3/2&0&0\end{pmatrix}.
\label{fo:relaxation-maximum}
\end{equation}
If $H$ is positive definite this maximum is strictly positive by
congruence and inertia.  This particular relaxation cannot prove an
opposite threshold sign.  If a Gram is singular, projection onto its
range, or a Moore--Penrose inverse, is required; positivity is not
asserted for every actual height.

The prescribed coefficients themselves give an additional exact
Gaussian relation.  Set
\begin{equation}
b_n=e^{-\sigma\mu+t\mu^2/4}
          e^{-(\sigma-t\mu/2)\rho_n},\qquad
a_n(Z)=b_n e^{\sqrt{t/2}\,Z\rho_n},
\label{fo:Gaussian-tilt}
\end{equation}
where $Z$ is a standard real Gaussian.  Completing the square shows
$\mathbb Ea_n(Z)=w_n$.  For the actual five-vector shape
$f_n=r(\rho_n)\cos\phi_n+s(\rho_n)\sin\phi_n$, let
$V(Z)=\sum_na_n(Z)f_n$ and $V=\mathbb EV(Z)=(e,u)$.  Then
\begin{align}
\mathcal K_{\Lambda,B}(V)
 &=\mathbb E\mathcal K_{\Lambda,B}(V(Z))
                  -\operatorname{tr}(\mathbb Q\mathcal C_t),
 \label{fo:covariance}\\
\mathcal C_t
 &=\sum_{n,m\le N}w_nw_m
       (e^{t\rho_n\rho_m/2}-1)f_nf_m^T
 =\sum_{k\ge1}\frac{(t/2)^k}{k!}V^{[k]}(V^{[k]})^T,\nonumber\\
V^{[k]}
 &=\left(X_k,Y_{k+1}+\epsilon X_{k+1},
                      X_{k+2},Y_{k+3},X_{k+4}\right)^T.\nonumber
\end{align}
The sums over $n,m$ are finite and the series converges absolutely,
so Gaussian integration and exponential expansion justify all
interchanges.  The same four kernels of \eqref{fo:four-channels}
occur in the covariance, multiplied inside every pair coefficient by
$e^{t\rho_n\rho_m/2}-1$.  The contraction of a positive covariance
with the indefinite $\mathbb Q$ has no automatic sign.  Only the mean
lower coordinates are constrained by the actual candidate; neither
each tilted state nor each shifted vector $V^{[k]}$ is a candidate.

For completeness the specified absolute covariance envelope has size
\begin{equation}
\mathcal E_t:=\sum_{n,m\le N}w_nw_m
 |e^{t\rho_n\rho_m/2}-1|\,\|f_n\|\|f_m\|
 \asymp t(Nw_N)^2.
\label{fo:covariance-envelope}
\end{equation}
Here is an elementary bound retaining the complete cutoff.  In bins
$j\le y_n=\log(N/n)<j+1$, there are at most $2Ne^{-j}$ indices,
$|\rho_n|\le y_n+O(N^{-1})$, and
$\|f_n\|\le C(1+y_n)^4$.  The inequality
$|e^z-1|\le|z|e^{|z|}$ combines the two heat exponents and its
covariance multiplier into $t(y_n+y_m+o(1))^2/4$.
After the bin counts the exponential factor is at most
\begin{equation}
\exp\{-(1/2-o(1))(j+k)+t(j+k+2+o(1))^2/4+O(1)\}.
\end{equation}
Since $j+k\le2\log N=\kappa/t+o(1)$ and $\kappa\le3/2$,
it is bounded by a constant times $e^{-(j+k)/16}$ for sufficiently
small time.  All fixed polynomial factors are summable, proving the
upper bound in \eqref{fo:covariance-envelope}.  For the lower bound
keep $Ne^{-2}\le n,m\le Ne^{-1}$.  This supplies a fixed positive
proportion of pairs, each with weight at least a fixed multiple of
$w_N^2$, $\rho_n\rho_m\ge0.81$, and
\begin{equation}
\|f_n\|^2\ge\cos^2\phi_n+\rho_n^6\sin^2\phi_n\ge0.9^6.
\end{equation}
Also $e^{t\rho_n\rho_m/2}-1\ge0.405t$.
This proves the lower bound for every carrier and common height.
Thus a triangle payment $\|\mathbb Q\|\mathcal E_t$ is exponentially
large, since $Nw_N\asymp e^{\mathfrak a/t}$.  This diagnoses that
specified absolute covariance method, not the size or sign of the
actual signed contraction in \eqref{fo:covariance}.

\section{A diagonal-annihilating extension and the near-pair theorem}
\label{fo:annihilation}

Define the complete finite-moment extension
\begin{equation}
\mathcal J^\dagger=\mathcal K
       +X_0(\Gamma X_4-3X_6+2\epsilon Y_6)-2Z_1Y_5.
\label{fo:Jdag}
\end{equation}
It agrees with $\mathcal K$ on the exact finite candidate, and
\begin{equation}
\mathcal K-\mathcal J^\dagger
 =X_0(3X_6-\Gamma X_4-2\epsilon Y_6)+2Z_1Y_5.
\label{fo:Jdag-difference}
\end{equation}
Two valid measured candidate payments are
\begin{align}
\Pi^\dagger_{\rm box}
 &=a_0|3X_6-\Gamma X_4-2\epsilon Y_6|+2a_1|Y_5|,
 \label{fo:Jdag-payment}\\
\Pi^\dagger_{\rm slab}
 &=\frac{\eta_N}{2}
 \left|3X_6-\Gamma X_4-2\epsilon Y_6+\frac{2a_{\rm c}}cY_5\right|
             +\frac{L\eta_N}{c}|Y_5|.\nonumber
\end{align}
The second follows by substituting
$Z_1=(a_{\rm c}X_0-e_1)/c$ into \eqref{fo:Jdag-difference}.
Either, or their minimum, bounds
$|\mathcal K-\mathcal J^\dagger|$.  Lemma~\ref{fo:moment-lemma}
applies to $j=5,6$, and
$\epsilon=O(t/x)$, $a_{\rm c}=O(x^{-1})$, $\Gamma=O(L)$ give
\begin{equation}
\Pi^\dagger_{\rm box}=O(L\eta_NNw_N)
                        =O(LN^{-\kappa/4})=o(1).
\label{fo:Jdag-payment-scale}
\end{equation}
Only the two genuine lower-jet tolerances are used.
The objects $X_6,Y_6,Y_5$ are exact finite logarithmic moments; they
are not identified with genuine fifth or sixth spatial jets, and no
such jet approximation is asserted.  The older dual family in
\eqref{fo:dual-definition} adds features of degree at most five to
the pair kernel and cannot identically remove its degree-six parts.
The new expression extends that feature family; it cancels the diagonal
rather than the entire kernel.

\begin{theorem}[Four exact diagonal-annihilating kernels]
\label{fo:Jdag-kernels}
For $a=\rho_n$, $b=\rho_m$, $h=a+b$, and $\delta=a-b$, one has
\begin{align}
\mathcal J^\dagger=\sum_{n,m\le N}w_nw_m\{&
D_c^\dagger\cos\Delta_{nm}+D_s^\dagger\sin\Delta_{nm}\nonumber\\
&+P_c^\dagger\cos\Sigma_{nm}+P_s^\dagger\sin\Sigma_{nm}\},
 \label{fo:Jdag-pairs}\\
D_c^\dagger&=\frac14h^2\delta^2(\Gamma-2h^2-\delta^2),
 &P_c^\dagger&=\frac14h^2\delta^2(\Gamma-h^2-2\delta^2),\nonumber\\
D_s^\dagger&=\frac\epsilon2(a-b)(a^5+b^5),
 &P_s^\dagger&=\frac\epsilon2(a-b)(a^5-b^5).\nonumber
\end{align}
All four kernels vanish individually on $n=m$.
\end{theorem}
\begin{proof}
The expansion of \eqref{fo:Jdag}, symmetrized in its full ordered
pair sum, has a summand
\begin{equation}
R\cos\phi_n\cos\phi_m+S\sin\phi_n\sin\phi_m
 +C_{nm}\cos\phi_n\sin\phi_m+C_{mn}\sin\phi_n\cos\phi_m,
\end{equation}
where
\begin{align}
R={}&-\Gamma a^2b^2+\frac32(a^2b^4+a^4b^2)
       +\frac\Gamma2(a^4+b^4)-\frac32(a^6+b^6),\nonumber\\
S={}&2a^3b^3-ab^5-a^5b,\qquad
C_{nm}=\epsilon b^5(b-a).\label{fo:Jdag-RSC}
\end{align}
For example the drift contribution is
$2\epsilon(X_0Y_6-X_1Y_5)$, whose unsymmetrized coefficient is
$2\epsilon(b^6-ab^5)\cos\phi_n\sin\phi_m$.
Product-to-sum therefore gives $(R+S)/2$, $(R-S)/2$,
$(C_{mn}-C_{nm})/2$, and $(C_{mn}+C_{nm})/2$ for the four
channels.  For the cosine factors,
\begin{equation}
R=(a^2-b^2)^2\{\Gamma/2-3(a^2+b^2)/2\},\qquad
S=-ab(a^2-b^2)^2.
\end{equation}
Use $a^2-b^2=h\delta$, $a^2+b^2=(h^2+\delta^2)/2$, and
$ab=(h^2-\delta^2)/4$ to obtain the claimed cosine kernels.
Also $C_{mn}=\epsilon a^5(a-b)$, so
\begin{equation}
C_{mn}-C_{nm}=\epsilon(a-b)(a^5+b^5),\qquad
C_{mn}+C_{nm}=\epsilon(a-b)(a^5-b^5).
\end{equation}
This verifies the difference-sine sign and both factors of $1/2$.
Every displayed kernel contains $\delta$; the product sine has
another factor because $a^5-b^5$ is divisible by $a-b$.
\end{proof}
The difference-sine coefficient is antisymmetric and its phase is
antisymmetric too.  Both cosine kernels and the product-sine kernel
are symmetric.  The ratio diagonal $(a,b)=(1,1)$ is now exactly zero.
For a product square only its diagonal divisor pair vanishes; its
off-diagonal divisor pairs remain.  The same true pair weight
\eqref{fo:pair-weight}, common height, and four block/complement classes
are retained.

\begin{theorem}[Uniform macroscopic near-pair bound]
\label{fo:near-pair}
Fix $0<\lambda<1$ and $1\le H\le N$.  For any subset $\mathcal E$
of the ordered pairs
\begin{equation}
\lambda N\le n,m\le N,\qquad |n-m|\le H,
\end{equation}
let
\begin{align}
\mathsf B_c&=\sum_{(n,m)\in\mathcal E}w_nw_m
                   (|D_c^\dagger|+|P_c^\dagger|),\nonumber\\
\mathsf B_s&=\sum_{(n,m)\in\mathcal E}w_nw_m
                   (|D_s^\dagger|+|P_s^\dagger|).
\end{align}
Uniformly for all sufficiently small positive time in
\eqref{fo:sector},
\begin{equation}
\mathsf B_c\le C_\lambda(1+|\Gamma|)w_N^2H^3/N,
\qquad
\mathsf B_s\le C_\lambda|\epsilon|w_N^2H^2.
\label{fo:near-pair-bounds}
\end{equation}
The product sine part alone is at most
$C_\lambda|\epsilon|w_N^2H^3/N$.
\end{theorem}
\begin{proof}
The centering estimate gives $|\rho_n|,|\rho_m|\le
1+\log(1/\lambda)$, and the mean value theorem gives
$|\delta|=|\log(n/m)|\le|n-m|/(\lambda N)$.
The exact weight-distance formula \eqref{fo:weight-distance}, now with
$0\le y_n\le\log(1/\lambda)$, gives $w_n,w_m\le C_\lambda w_N$.
Writing $k=|n-m|$, Theorem~\ref{fo:Jdag-kernels} yields
\begin{equation}
|D_c^\dagger|+|P_c^\dagger|
 \le C_\lambda(1+|\Gamma|)k^2/N^2,
\quad |D_s^\dagger|\le C_\lambda|\epsilon|k/N,
\quad |P_s^\dagger|\le C_\lambda|\epsilon|k^2/N^2.
\end{equation}
For each positive integer gap $k$, at most $2N$ ordered pairs have
that gap.  The gap-zero contribution is exactly zero.  Consequently
the squared-gap sum is at most
$(2/N)\sum_{k\le H}k^2\le CH^3/N$, and the linear-gap sum is at
most $2\sum_{k\le H}k\le CH^2$.  Multiplication by the weight
bound proves the theorem, using $H^3/N\le H^2$ for the total sine
bound.  This is an absolute bound for the stated band, requiring no
phase cancellation.
\end{proof}
Since
\begin{equation}
(Nw_N)^2\asymp N^{1-\kappa/4},\qquad x\asymp N^2,
\qquad \epsilon=O(t/N^2),
\end{equation}
the choice $H=\lceil\sqrt N\rceil$ gives
\begin{equation}
\mathsf B_c=O(LN^{-1/2-\kappa/4}),\qquad
\mathsf B_s=O(tN^{-2-\kappa/4}),
\label{fo:near-pair-scales}
\end{equation}
and the product sine part is $O(tN^{-5/2-\kappa/4})$.
Both are $o(N^{-\kappa/4})$.  Taking $\lambda=1/2$ covers the
macroscopic block.  Taking $\lambda=1/4$ also covers all near pairs
crossing its boundary for this choice of $H$ and sufficiently large
$N$.  Low-index near pairs outside this range remain exact.

Let $\mathcal J^\dagger_{\rm near}$ be precisely the part of
\eqref{fo:Jdag-pairs} with $\min(n,m)\ge N/4$ and
$|n-m|\le H$, and put
$\mathcal J^\dagger_{\rm rem}=\mathcal J^\dagger-
\mathcal J^\dagger_{\rm near}$.  An actual-phase bound
$\mathcal J^\dagger_{\rm rem}\le U_{\rm rem}$ would give
\begin{equation}
\mathcal K\le U_{\rm rem}+\mathsf B_c+\mathsf B_s+\Pi^\dagger.
\end{equation}
The fully paid sufficient ordinary-double exclusion is therefore
\begin{equation}
U_{\rm rem}+\mathsf B_c+\mathsf B_s+\Pi^\dagger
                    +\widehat\Delta_{\rm count}/(4c^6)<0.
\label{fo:final-paid-target}
\end{equation}
Any subsequent signed transformation must add its proved error.
An alternative representation of the same sum is not an additional
disjoint contribution.  The theorem pays a narrow band; it supplies no
sign for the retained separated primitive and product terms.

Finally, the algebraic antilinear map
$\mathcal R:M_j\mapsto(-1)^j\overline{M_j}$ fixes $X_{2j}$ and
$Y_{2j+1}$ and hence fixes $\mathcal K$ for the same frozen
$\Gamma$.  At fixed physical drift it gives
\begin{equation}
Z_1\mapsto Z_1-2\epsilon X_1,\qquad
\mathcal J^\dagger(\mathcal RM;\epsilon)
 =\mathcal J^\dagger(M;-\epsilon)
 =\mathcal J^\dagger(M;\epsilon)
       -4\epsilon X_0Y_6+4\epsilon X_1Y_5.
\label{fo:reflection-parity}
\end{equation}
The full moment envelope pays this fixed-drift discrepancy by
$4|\epsilon|(B_0B_6+B_1B_5)=O(tN^{-1-\kappa/4})$.
Its relation to the prescribed stationary transformation still requires
the transformed cutoff and boundary terms; the parity calculation alone
cannot change the sign of the threshold quadratic.

\section{Common-factor smoothing of the difference channel}
\label{em:section}
Let $D_c,D_s$ be the two difference polynomials of a frozen paid quadratic.
For both the original five-feature dual family and the
diagonal-annihilating extension they have total degree at most six.
Let $\mathfrak M$ bound the absolute values of their polynomial
coefficients. For example the original family admits the scale
$(1+|\epsilon|)^2(1+|\Gamma|+\|\Lambda\|+\|B\|)$.
All constants below retain this scale.

Every ordered pair is uniquely $n=ja,m=jb$ with $(a,b)=1$.
Write $J_{a,b}=\lfloor N/\max(a,b)\rfloor$, and define
\begin{align}
 W_{a,b}&=(ab)^{-\sigma}
       \exp\{t(\log^2a+\log^2b)/4\},&
 \beta_{a,b}&=2\sigma-\tfrac t2\log(ab),\notag\\
 H_{a,b,\chi}(z)&=D_\chi(z+\log a-\mu,z+\log b-\mu),&
 f_{a,b,\chi}(u)&=W_{a,b}u^{-\beta_{a,b}}
             e^{t\log^2u/2}H_{a,b,\chi}(\log u).
 \label{em:coefficient}
\end{align}
Then $f_{a,b,\chi}(j)=w(ja)w(jb)D_\chi(\rho_{ja},\rho_{jb})$,
and the exact difference channel is
\[
\sum_{\substack{(a,b)=1\\\max(a,b)\le N}}
\left\{\cos(T\log(a/b))\sum_{j=1}^{J_{a,b}}f_{a,b,c}(j)
      +\sin(T\log(a/b))\sum_{j=1}^{J_{a,b}}f_{a,b,s}(j)\right\}.
\]
The phase is independent of $j$ in this channel.

Let $K=\lfloor N/2\rfloor$, $\mathcal B=\{K+1,\ldots,N\}$,
$\mathcal C=\{1,\ldots,K\}$, and $k_a=\lfloor K/a\rfloor$,
$k_b=\lfloor K/b\rfloor$. The following literal integer intervals
partition $1\le j\le J_{a,b}$:
\[
\begin{array}{c|cc}
\text{class}&\ell&v\\ \hline
\mathcal C\times\mathcal C&1&\min(J_{a,b},k_a,k_b)\\
\mathcal B\times\mathcal B&\max(1,k_a+1,k_b+1)&J_{a,b}\\
\mathcal B\times\mathcal C&\max(1,k_a+1)&\min(J_{a,b},k_b)\\
\mathcal C\times\mathcal B&\max(1,k_b+1)&\min(J_{a,b},k_a)
\end{array}
\]
Rows with $\ell>v$ are empty. For an integer $J_*\ge1$, retain all
$j\le J_*$ exactly and transform only $p=\max(\ell,J_*+1)\le j\le v=q$.

\begin{lemma}[Uniform centered coefficient moments]
\label{em:moment-bound}
For every fixed integer $d\ge0$, uniformly for real $1\le u\le N$,
\[
\sum_{a\le N/u}w(ua)(1+|\log(ua)-\mu|)^d
                      \le C_d (N/u)w_N .
\]
\end{lemma}
\begin{proof}
This is \eqref{fo:real-index-moments}; the logarithmic bin proof in
Lemma~\ref{fo:moment-lemma} applies for real $u$ and includes the final
low-index bin. We state it here to emphasize the real-index estimate
needed inside the remainder integrals.
\end{proof}

Fix an integer order $r\ge1$. For an integer interval $p\le q$, define
\[
\mathsf E_r[f;p,q]=\int_p^q f(u)\dd u+\frac{f(p)+f(q)}2
 +\sum_{h=1}^r\frac{B_{2h}}{(2h)!}
                      \{f^{(2h-1)}(q)-f^{(2h-1)}(p)\}.
\]
Euler--Maclaurin and the Bernoulli Fourier series \cite{dlmfem,dlmfbern}
give exactly
\begin{equation}
 \sum_{j=p}^q f(j)=\mathsf E_r[f;p,q]
 -\frac1{(2r)!}\int_p^q\widetilde B_{2r}(u)f^{(2r)}(u)\dd u,
 \qquad
 |R_r|\le\frac{2\zeta(2r)}{(2\pi)^{2r}}\int_p^q|f^{(2r)}(u)|\dd u .
 \label{em:formula}
\end{equation}
Singleton intervals have zero remainder. All derivative corrections are
explicit: with $P_0=H_{a,b,\chi}$ and
$P_{k+1}(z)=P_k'(z)+(tz-\beta_{a,b}-k)P_k(z)$,
\[
f_{a,b,\chi}^{(k)}(u)=
W_{a,b}u^{-\beta_{a,b}-k}e^{t\log^2u/2}P_k(\log u).
\]
The recurrence follows by $d/du=u^{-1}d/dz$.
If $H(z)=\sum_{m=0}^6h_mz^m$ and $\lambda=1-\beta_{a,b}$,
the main integral is
\[
W_{a,b}\sum_{m=0}^6h_m\partial_\lambda^m
\left\{\sqrt{\frac\pi{2t}}e^{-\lambda^2/(2t)}
\left[\erfi\left(\frac{tz+\lambda}{\sqrt{2t}}\right)\right]
                          _{\log p}^{\log q}\right\}.
\]

\begin{theorem}[Global paid difference transformation]
\label{em:remainder}
Replace every large-common-factor interval by $\mathsf E_r$, keeping
all floors, all four pair classes, and all $j\le J_*$ exact.
The complete difference-channel error is bounded by
\[
 C_r\mathfrak M(Nw_N)^2J_*^{-2r-1}.
\]
For the density-strengthened coefficient $\Gamma=O(L)$ and bounded
other dual parameters, $r=1$, $J_*=\lceil N^{2/5}\rceil$ give
$O(LN^{-1/5-\kappa/4})=o(N^{-\kappa/4})$.
\end{theorem}
\begin{proof}
The recurrence and bounded physical quantities $t\log u,\beta_{a,b}$
show, for each fixed $k$,
\[
|f_{a,b,\chi}^{(k)}(u)|\le C_k\mathfrak M u^{-k}w(ua)w(ub)
(1+|\rho(ua)|)^6(1+|\rho(ub)|)^6 .
\]
Differentiation increases only a fixed polynomial degree; its coefficients
remain uniformly bounded because $t\le1/20$ and $t\log u$ is bounded.
At a fixed $u$, drop coprimality for this positive error envelope and apply
Lemma~\ref{em:moment-bound} twice. This contributes at most
$C_r\mathfrak M (Nw_N)^2u^{-2r-2}$.
The four intervals are disjoint for each pair, so their remainder
integrals do not multiply the coverage of $u$; the two kernels add
only a fixed constant. Integrating for $u\ge J_*$ proves the bound.
Finally $(Nw_N)^2\asymp N^{1-\kappa/4}$ gives the stated power.
\end{proof}

This theorem does not smooth the product phase by absolute derivatives.
That phase contains $u^{2iT}$; inserting it into the same derivative
envelope introduces $(1+|T|)^{2r}$. Since $T\asymp N^2$, the resulting
upper envelope grows even when $J_*$ is comparable to $N$, for each
fixed $r\ge1$. This is a loss of that specific envelope, not a lower
bound for the true oscillatory remainder. The next section keeps the
product phase inside the oscillatory integrals.

\section{Oscillation-preserving product transformations}
\label{tr:product-section}

All data in this section are frozen at one physical center.  Derivatives
are derivatives of an auxiliary summation variable, never raw spatial
jet derivatives.  Write
\[
 w(u)=u^{-\sigma}\exp\!\left(\frac t4\log^2u\right),\qquad
 \rho_u=\log u-\mu,\qquad \phi_u=\theta_t+T\log u,\qquad
 P=\frac{T}{2\pi}.
\]
We use the sector $L=\kappa/t$, $1\le\kappa\le3/2$, with the literal
natural cutoff $N=\lfloor\sqrt{e^L+t/16}\rfloor$.  The prescribed physical
center has, uniformly as $t\downarrow0$,
\begin{equation}
\label{tr:physical-scales}
 T\asymp N^2,\quad P/N^2\longrightarrow1,\quad \mu=\log N+O(N^{-1}),\quad
 c_\sigma:=\sigma-\frac t2\log N=\frac12+o(1),\quad
 t\log N=\frac\kappa2+o(1),\quad
 Nw_N\asymp N^{1/2-\kappa/8},
\end{equation}
where $w_N=w(N)$.  All asymptotic constants below are uniform for
sufficiently small $t$ on this closed $\kappa$ interval.  Exact identities
keep their stated data and floors.

For small enough time we use the numerical reserves
\begin{equation}
\label{tr:reserves}
 \frac7{16}\le c_\sigma\le\frac9{16},\qquad
 \frac7{16}\le t\log N\le\frac{13}{16},\qquad
 |\mu-\log N|\le1.
\end{equation}
The centered sublattice moment bound proved in
Lemma~\ref{em:moment-bound} gives, for every fixed nonnegative integer $d$,
\begin{equation}
\label{tr:weighted-moments}
 \sum_{u\le N/q}w(qu)(1+|\rho_{qu}|)^d
 \le C_d\frac Nq w_N\qquad(1\le q\le N).
\end{equation}
Sums over real upper limits mean sums over integers not exceeding that limit.

Consider a product pair kernel
\[
 w(n)w(m)\{P_c(n,m)\cos(2\theta_t+T\log(nm))
             +P_s(n,m)\sin(2\theta_t+T\log(nm))\},
\]
where $P_c,P_s$ are real polynomials in $(\rho_n,\rho_m)$ of total
 degree at most six.  Let $\mathfrak M\ge1$ bound their coefficients,
including their frozen drift factors.  The original paid dual and the
extended diagonal-annihilating dual both have this degree.  Define
\begin{equation}
\label{tr:product-amplitude}
 A_{a,b}(u)=w(ua)w(ub)\{P_c(ua,ub)-iP_s(ua,ub)\}.
\end{equation}
The sign is fixed by
$\operatorname{Re}[(P_c-iP_s)e^{i\Phi}]=P_c\cos\Phi+P_s\sin\Phi$.
Unique common-factor decomposition gives the exact product channel
\begin{equation}
\label{tr:product-cf}
 \mathcal P=\operatorname{Re}\left[
 e^{2i\theta_t}\sum_{\substack{(a,b)=1\\\max(a,b)\le N}}
 e^{iT\log(ab)}
 \sum_{j=1}^{\lfloor N/\max(a,b)\rfloor}
 A_{a,b}(j)e^{2iT\log j}\right].
\end{equation}

The literal block/core partition is retained before transformation.
Put $K_N=\lfloor N/2\rfloor$, $k_a=\lfloor K_N/a\rfloor$,
$k_b=\lfloor K_N/b\rfloor$, and $J_{a,b}=\lfloor N/\max(a,b)\rfloor$.
The four integer intervals are
\[
\begin{array}{c|cc}
 \text{class}&\text{left endpoint}&\text{right endpoint}\\\hline
 \mathcal C\times\mathcal C&1&\min(J_{a,b},k_a,k_b)\\
 \mathcal B\times\mathcal B&\max(1,k_a+1,k_b+1)&J_{a,b}\\
 \mathcal B\times\mathcal C&\max(1,k_a+1)&\min(J_{a,b},k_b)\\
 \mathcal C\times\mathcal B&\max(1,k_b+1)&\min(J_{a,b},k_a).
\end{array}
\]
Only empty intervals are discarded.  Intersect every remaining interval
with the disjoint integer cells $[U,2U-1]$, $U=2^h$, $h\ge0$.
Their resulting endpoints $[p,q]$ satisfy
$U\le p\le q<2U$ and $qa,qb\le N$.  Singletons are kept exactly.

\subsection{A finite interval theorem with endpoint transitions}

Let $\tau>0$ denote an oscillation parameter, distinct from the heat
time $t$.  For a nonsingleton integer interval $[p,q]$ in one dyadic
cell, put
\[
 \Sigma_\tau[A;p,q]=\sum_{j=p}^qA(j)e^{2i\tau\log j},\qquad
 \psi_k(u)=2\tau\log u-2\pi ku,\qquad
 I_k=\int_p^qA(u)e^{i\psi_k(u)}\,du.
\]
Define the finite expanded stationary band
\begin{equation}
\label{tr:stationary-band}
 \mathcal K=\left\{k\in\mathbb Z_{>0}:
       \frac{\tau}{2\pi q}\le k\le\frac{2\tau}{\pi p}\right\}.
\end{equation}
For $k\in\mathcal K$, $u_k=\tau/(\pi k)\in[p/2,2q]$.
For excluded modes write $V_k(u)=2\tau/u-2\pi k$ and
\[
 S_m(u)=\sum_{k\notin\mathcal K}V_k(u)^{-m}\quad(m=2,3),\qquad
 S_1(u)=\lim_{R\to\infty}
 \sum_{\substack{|k|\le R\\k\notin\mathcal K}}V_k(u)^{-1}.
\]
The last limit pairs positive and negative modes.  Its summands beyond
any fixed finite set have canceling leading $1/k$ terms.  Define
\begin{equation}
\label{tr:outside-boundary}
 B_{\mathrm{out}}[A;p,q]=
 \left[e^{2i\tau\log u}
 \left\{\frac{A(u)}i S_1(u)+A'(u)S_2(u)
             +\frac{2\tau A(u)}{u^2}S_3(u)\right\}\right]_p^q.
\end{equation}
These are explicit endpoint functions.  With $z=\tau/(\pi u)$ and
$\mathcal K$ fixed, one can equivalently use
\begin{equation}
\label{tr:cotangent}
 C_1(z)=\pi\cot(\pi z)-\sum_{k\in\mathcal K}\frac1{z-k},\qquad
 C_2=-C_1',\qquad C_3=\tfrac12 C_1'',\qquad
 S_m=(2\pi)^{-m}C_m.
\end{equation}
Any integer $z$ occurring for $p\le u\le q$ belongs to $\mathcal K$,
so the apparent poles are removable.  The cotangent partial-fraction
identity \cite{dlmfcot} gives the same paired series as above.

For the retained modes introduce the exact quadratic coordinate
\[
 \xi(v)=\operatorname{sgn}(v-1)\sqrt{2(v-1-\log v)},\qquad
 \xi(1)=0,
\]
and let $v=\mathcal V(\xi)$ be its smooth increasing inverse.  Set
\[
 a_k=\xi(p/u_k),\quad b_k=\xi(q/u_k),\qquad
 g_k(\xi)=u_kA(u_k\mathcal V(\xi))\mathcal V'(\xi).
\]
For smooth $g$ define regular operators
\[
 \mathcal Bg(\xi)=\frac{g(\xi)-g(0)}\xi,
 \quad \mathcal Bg(0)=g'(0),\qquad
 \mathcal Dg=(\mathcal Bg)',\qquad
 \mathcal F_\tau(a,b)=\int_a^be^{-i\tau\xi^2}\,d\xi.
\]
The finite retained-mode expression is
\begin{align}
\label{tr:retained-Fresnel}
 F_k[A;p,q]=e^{i\psi_k(u_k)}\bigg\{&
 \left(g_k(0)+\frac{g_k''(0)}{4i\tau}\right)
       \mathcal F_\tau(a_k,b_k)\notag\\
 &-\frac{[\mathcal Bg_k(\xi)e^{-i\tau\xi^2}]_{a_k}^{b_k}}{2i\tau}
 -\frac{[\mathcal B\mathcal Dg_k(\xi)e^{-i\tau\xi^2}]_{a_k}^{b_k}}
              {(2i\tau)^2}\bigg\}.
\end{align}
Here
\[
 g_k(0)=u_kA(u_k),\qquad
 g_k''(0)=u_k\{u_k^2A''(u_k)+2u_kA'(u_k)+A(u_k)/6\},
\]
and the entire-function representation is
\[
 \mathcal F_\tau(a,b)=
 \frac{\sqrt\pi e^{-i\pi/4}}{2\sqrt\tau}
 \left[\operatorname{erf}(e^{i\pi/4}\sqrt\tau\,\xi)\right]_a^b.
\]
In particular the curvature sign is negative and all regular quotients
remain defined when an endpoint is stationary.

\begin{theorem}[Finite product Poisson--Fresnel transformation]
\label{tr:finite-product}
Suppose $\tau\ge U$ and $A$ is $C^4$ on $[U/2,4U]$, with
\[
 |A^{(h)}(u)|\le C_h\mathfrak M U^{-h}W_U
       \quad(0\le h\le4).
\]
Define
\begin{equation}
\label{tr:finite-main}
 \mathsf T_\tau[A;p,q]=
 \tfrac12\{A(p)e^{2i\tau\log p}+A(q)e^{2i\tau\log q}\}
 +B_{\mathrm{out}}[A;p,q]
 +\sum_{k\in\mathcal K}F_k[A;p,q].
\end{equation}
Then
\[
 |\Sigma_\tau[A;p,q]-\mathsf T_\tau[A;p,q]|
 \le C\mathfrak M W_U/\tau.
\]
The constant depends only on the derivative-envelope constants.  No
separation of a stationary point from an endpoint is assumed.  For a
singleton define $\mathsf T_\tau$ to be its exact summand.
\end{theorem}
\begin{proof}
Periodize the finitely many unit intervals by
$G(v)=\sum_{j=p}^{q-1}A(j+v)e^{2i\tau\log(j+v)}$, $0<v<1$.
Its $k$th Fourier coefficient is $I_k$.  The symmetric Fourier series
at zero converges to the average of the two one-sided limits.  Thus
\begin{equation}
\label{tr:finite-poisson}
 \Sigma_\tau[A;p,q]=
 \tfrac12\{A(p)e^{2i\tau\log p}+A(q)e^{2i\tau\log q}\}
 +\lim_{R\to\infty}\sum_{|k|\le R}I_k.
\end{equation}
This elementary piecewise-smooth Fourier convergence is the only
Fourier input \cite{dlmffourier}.

Outside $\mathcal K$, $V_k$ never vanishes.  For $k\le0$ its
magnitude is at least $2\tau/q+2\pi|k|$; for
$0<k<\tau/(2\pi q)$ it is at least $\tau/q$; and for
$k>2\tau/(\pi p)$ it is at least $\pi k$.  Counting the low positive
range and integrating the two tails gives, uniformly on $[p,q]$,
\begin{equation}
\label{tr:inverse-phase-bounds}
 \sum_{k\notin\mathcal K}|V_k|^{-2}\le CU/\tau,
 \qquad \sup_{k\notin\mathcal K}|V_k|^{-1}\le CU/\tau.
\end{equation}
Write $\mathcal L_kA=(A/V_k)'$.  Two exact integrations by parts give
\[
 I_k=\left[e^{i\psi_k}
       \left\{\frac A{iV_k}+\frac{\mathcal L_kA}{V_k}\right\}
       \right]_p^q
       -\int_p^qe^{i\psi_k}\mathcal L_k^2A\,du,
\]
where
\[
 \mathcal L_k^2A=
 \frac{A''}{V_k^2}-\frac{3A'V_k'}{V_k^3}
 -\frac{AV_k''}{V_k^3}+\frac{3A(V_k')^2}{V_k^4}.
\]
At integer endpoints $e^{-2\pi iku}=1$, so the summed boundaries are
exactly \eqref{tr:outside-boundary}.  Since
$|V_k'|\le C\tau/U^2$ and $|V_k''|\le C\tau/U^3$,
\eqref{tr:inverse-phase-bounds} and the amplitude bounds imply
\[
 \sum_{k\notin\mathcal K}|\mathcal L_k^2A(u)|
 \le C\mathfrak M W_U/(U\tau).
\]
Multiplication by the interval length $\le U$ pays the entire excluded
integral tail by $C\mathfrak M W_U/\tau$.

For a retained mode the coordinate is exact:
$\psi_k(u)=\psi_k(u_k)-\tau\xi^2$, and
$I_k=e^{i\psi_k(u_k)}\int_{a_k}^{b_k}g_k e^{-i\tau\xi^2}\,d\xi$.
Expanding the inverse at zero gives
$\mathcal V(\xi)=1+\xi+\xi^2/3+\xi^3/36+O(\xi^4)$ and hence the
stated formulas for $g_k(0)$ and $g_k''(0)$.  Direct integration by
parts gives
\[
 \int_a^bge^{-i\tau\xi^2}\,d\xi
 =g(0)\mathcal F_\tau(a,b)
 -\frac{[\mathcal Bg\,e^{-i\tau\xi^2}]_a^b}{2i\tau}
 +\frac1{2i\tau}\int_a^b\mathcal Dg\,e^{-i\tau\xi^2}\,d\xi.
\]
Applying this twice yields \eqref{tr:retained-Fresnel} and the exact
remainder
\[
 I_k-F_k=
 \frac{e^{i\psi_k(u_k)}}{(2i\tau)^2}
 \int_{a_k}^{b_k}\mathcal D^2g_k(\xi)e^{-i\tau\xi^2}\,d\xi.
\]
The regular quotient satisfies
\[
 \mathcal Dg(\xi)=\int_0^1s\,g''(s\xi)\,ds,\qquad
 \mathcal D^2g(\xi)=\int_0^1\!\int_0^1r s^3
              g^{(4)}(sr\xi)\,ds\,dr,
\]
so $\|\mathcal D^2g\|_\infty\le\|g^{(4)}\|_\infty/8$.
All required ratios $u/u_k$ lie in $[1/4,4]$; the inverse and its first
five derivatives are bounded on this compact interval.  The amplitude
hypothesis therefore gives $\|g_k^{(4)}\|_\infty\le
C\mathfrak M UW_U$, and the $\xi$ interval has bounded length.
The per-mode error is $C\mathfrak M UW_U/\tau^2$.  There are at most
$C\tau/U$ retained modes, so their total error is
$C\mathfrak M W_U/\tau$.  Combining both errors with
\eqref{tr:finite-poisson} proves the theorem.
\end{proof}

\subsection{Global product payments}

\begin{theorem}[Complete product-channel transformation]
\label{tr:global-product}
Let $\widetilde{\mathcal P}^{\mathrm{CF}}$ be obtained from
\eqref{tr:product-cf} by replacing every literal masked dyadic interval
with \eqref{tr:finite-main}, with $\tau=T$, before taking the real
part.  Then
\[
 |\mathcal P-\widetilde{\mathcal P}^{\mathrm{CF}}|
 \le C\mathfrak M (Nw_N)^2/T
 =O(\mathfrak M N^{-1-\kappa/4}).
\]
Alternatively, hold $m$ fixed, apply the same theorem to the original
$n$ index with $\tau=T/2$ and amplitude
$A_m(u)=w(u)w_m\{P_c(u,m)-iP_s(u,m)\}$, retaining the literal block
and core intervals for both indices.  Provided the physical relation
\begin{equation}
\label{tr:sharper-sigma}
 \sigma-\frac t4\log N=s_\kappa+O(t/N),
 \qquad s_\kappa=\frac12+\frac\kappa8,
\end{equation}
holds, its full transformed product satisfies
\[
 |\mathcal P-\widetilde{\mathcal P}^{\mathrm{orig}}|
 \le C\mathfrak M L^6(Nw_N)/T
 =O(\mathfrak M L^6N^{-1-s_\kappa}).
\]
Both formulas transform the entire product channel, including common
factor $j=1$.  All singleton, mixed-class, and endpoint terms are retained.
\end{theorem}
\begin{proof}
For the common-factor amplitude use
\[
 W_{a,b,U}=w(Ua)w(Ub)
 (1+|\rho_{Ua}|)^6(1+|\rho_{Ub}|)^6.
\]
The logarithmic weight derivative $-\sigma+(t/2)\log v$ and its next
logarithmic derivative $t/2$ are bounded on $1/2\le v\le4N$.
Changing an index by a factor in $[1/2,4]$ changes its weight by a
bounded factor and its centered log by a bounded amount.  Four amplitude
derivatives therefore satisfy the hypothesis of
Theorem~\ref{tr:finite-product} with this envelope.  The real-index
extension beyond the cutoff is only an amplitude device; it introduces
no additional heat approximation.
Equation~\eqref{tr:weighted-moments}, with $q=U$, gives
\[
 \sum_{a\le N/U}w(Ua)(1+|\rho_{Ua}|)^6
 \le C(N/U)w_N.
\]
For each dyadic cell and coprime pair there are at most four mask
intervals.  Dropping coprimality only in the positive error envelope,
\[
 \sum_{a,b\le N/U}W_{a,b,U}\le C(Nw_N)^2/U^2.
\]
Now sum the finite theorem and use $\sum_{h\ge0}2^{-2h}=4/3$.
The first bound follows from \eqref{tr:physical-scales}.

For the original-index version an envelope is
$W_{m,U}=w(U)w_m(1+|\rho_U|)^6(1+|\rho_m|)^6$.
Equation~\eqref{tr:sharper-sigma} implies $w(U)\le CU^{-s_\kappa}$
for $1\le U\le N$: the adverse exponent error times $\log U$ is at
most $O(t\log N/N)=O(N^{-1})$.  Since $|\rho_U|\le CL$,
\[
 \sum_{U=2^h\le N}w(U)(1+|\rho_U|)^6\le CL^6.
\]
The complete $m$ moment is $O(Nw_N)$ by
Equation~\eqref{tr:weighted-moments}.  Sum the finite theorem with
$\tau=T/2$ to obtain the second bound.  For the prescribed center,
\eqref{tr:sharper-sigma} follows from
$\operatorname{Re}\alpha-L/2=O(x^{-2})$ and
$\log N=L/2+O(N^{-1})$.
\end{proof}

With the density-strengthened dual $\mathfrak M=O(L)$, both product
remainders are $o(N^{-\kappa/4})$.  This comparison uses the available
physical upper-budget scale.  It asserts no positive lower bound for the
measured physical error, and the actual candidate and physical payments
remain necessary.

\section{The coupled primitive frontier}
\label{tr:primitive-section}

Let $v(\rho)$ be a complex vector polynomial of fixed dimension and degree
at most $D$.  Choose $A_v\ge1$ bounding its coefficient norms, and let
$Q$ be a real symmetric matrix.  Define
\[
 z_n=w(n)e^{i\phi_n}v(\rho_n),\qquad y_n=\operatorname{Re}z_n.
\]
The original five-feature dual has
\[
 v(\rho)=r(\rho)-is(\rho),\quad
 r=(1,\epsilon\rho,\rho^2,0,\rho^4)^T,\quad
 s=(0,\rho,0,\rho^3,0)^T,
\]
so $D=4$ and $A_v\le C(1+|\epsilon|)$.  The extended dual uses $D=6$.
For any such vector,
\begin{equation}
\label{tr:complete-pair}
 y_n^TQy_m=\tfrac12\operatorname{Re}
 \{z_n^TQ\overline{z_m}+z_n^TQz_m\}.
\end{equation}
These are the complete ratio and product terms, with both sine
orientations fixed by the actual complex features.

\begin{proposition}[Exact coupled M\"obius frontier]
\label{tr:mobius}
For $1\le J\le N$ set
\[
 \mathcal F_J=\sum_{\substack{n,m\le N\\\gcd(n,m)\le J}}y_n^TQy_m,
 \quad
 c_J(q)=\sum_{\substack{j\mid q\\j\le J}}\mu_{\mathrm{Mob}}(q/j),
 \quad V_q=\sum_{u=1}^{\lfloor N/q\rfloor}z_{qu}.
\]
Then
\begin{equation}
\label{tr:mobius-formula}
 \mathcal F_J=\sum_{q\le N}c_J(q)
 (\operatorname{Re}V_q)^TQ(\operatorname{Re}V_q)
 =\tfrac12\operatorname{Re}\sum_{q\le N}c_J(q)
 \{V_q^TQ\overline{V_q}+V_q^TQV_q\}.
\end{equation}
Moreover $|c_J(q)|\le\tau(q)$, $c_J(1)=1$, and
$c_J(q)=0$ for $1<q\le J$.  For $J=1$ the coefficients are
$\mu_{\mathrm{Mob}}(q)$; for $J=N$ only $q=1$ survives.
\end{proposition}
\begin{proof}
For any positive integer $g$,
\[
 \sum_{q\mid g}c_J(q)=
 \sum_{\substack{j\mid g\\j\le J}}
 \sum_{d\mid g/j}\mu_{\mathrm{Mob}}(d)
 =\mathbf1_{\{g\le J\}}.
\]
Insert this identity at $g=\gcd(n,m)$, interchange finite sums, and sum
over $n=qu,m=qv$.  Formula~\eqref{tr:complete-pair} gives the second
equality.  The coefficient checks follow immediately from the same
ordinary divisor identity.
\end{proof}

The $q>J$ terms generally survive; they cannot be truncated at $J$.
With $K_N=\lfloor N/2\rfloor$, let
\[
 m_q=\lfloor N/q\rfloor,\quad k_q=\lfloor K_N/q\rfloor,\quad
 V_q^{\mathcal C}=\sum_{u=1}^{k_q}z_{qu},\quad
 V_q^{\mathcal B}=\sum_{u=k_q+1}^{m_q}z_{qu}.
\]
Expanding $V_q=V_q^{\mathcal C}+V_q^{\mathcal B}$ in
\eqref{tr:mobius-formula} retains both mixed orientations and both
same-class terms.  The genuine paid candidate coordinates constrain the
complete vector $\operatorname{Re}V_1$, not each sublattice vector.
No separate candidate vanishing condition for $q>1$ is implied.

\subsection{An endpoint-paid sublattice transformation}

For one nonempty literal integer interval $p\le u\le r$ use
$a_0=p-1/2$, $b_0=r+1/2$.  For the core and block these are respectively
$[1/2,k_q+1/2]$ and $[k_q+1/2,m_q+1/2]$; empty sums give zero.
Define
\[
 a_q(u)=w(qu)v(\rho_{qu}),\qquad
 \Phi_q(u)=\theta_t+T\log(qu),\qquad
 I_{q,k}[a_0,b_0]=\int_{a_0}^{b_0}
 a_q(u)e^{i\Phi_q(u)-2\pi iku}\,du.
\]
Choose an integer $M\ge4P$.  At a half integer $z\ge1/2$ let
$\omega=P/z$ and
\[
 S_\ell(\omega,M)=\sum_{|k|>M}^{\mathrm{sym}}
       \frac{(-1)^k}{(\omega-k)^\ell},\qquad \ell=1,2,3.
\]
There is no pole, since $|\omega|\le2P\le M/2$.  The first series
converges absolutely after pairing $k,-k$:
$(\omega-k)^{-1}+(\omega+k)^{-1}=2\omega/(\omega^2-k^2)$.
The other two converge absolutely without pairing.  These endpoint
coefficients have a fixed-integral representation.  For $z>0$ put
\[
 F(z)=\int_0^1\frac{s^{z-1}}{1+s}\,ds
     =\sum_{n\ge0}\frac{(-1)^n}{n+z}.
\]
Integrating the finite geometric expansion proves the equality.  Then
\begin{equation}
\label{tr:tail-functions}
 S_1(\omega,M)=(-1)^{M+1}
 \{F(M+1+\omega)-F(M+1-\omega)\},\qquad
 S_2=-\partial_\omega S_1,\quad S_3=\tfrac12\partial_\omega^2S_1.
\end{equation}
Define the exact endpoint term and finite transformed vector by
\begin{align}
\label{tr:sublattice-endpoint}
 E_{q,M}(z)=e^{i\Phi_q(z)}\bigg\{&
 \frac{a_q(z)}{2\pi i}S_1(\omega,M)
 +\frac{a_q'(z)}{(2\pi)^2}S_2(\omega,M)\notag\\
 &+\frac{T a_q(z)}{z^2(2\pi)^3}S_3(\omega,M)\bigg\},\\
\label{tr:sublattice-main}
 \widetilde V_q[a_0,b_0]&=
 \sum_{|k|\le M}I_{q,k}[a_0,b_0]
 +E_{q,M}(b_0)-E_{q,M}(a_0).
\end{align}
All retained resonant integrals are exact.  In particular no
full-Gaussian replacement is part of this formula.

\begin{theorem}[Uniform sublattice payment]
\label{tr:sublattice-error}
Under \eqref{tr:reserves}, for every such literal interval,
\[
 \left\|\sum_{u=p}^rz_{qu}-\widetilde V_q[a_0,b_0]\right\|
 \le \frac{C_DA_v}{M}w(q)(1+|\rho_q|)^D.
\]
Constants may depend on $D$, dimension, and the fixed norm convention,
but not on $q,J,N,M$, or the frozen physical center.
\end{theorem}
\begin{proof}
Periodize the compactly supported function
$\mathbf1_{[a_0,b_0]}a_qe^{i\Phi_q}$.  Its periodization is piecewise
smooth and continuous at the integer evaluation point, because both
support endpoints are half integers.  Fourier convergence gives exactly
\[
 \sum_{u=p}^rz_{qu}=\lim_{R\to\infty}
                       \sum_{|k|\le R}I_{q,k}[a_0,b_0].
\]
There is no half weight on an integer summand.
For $|k|>M$ set $g_k=T/u-2\pi k$ and
$\mathcal D_k a=(a/(ig_k))'$.  Twice integrating by parts yields
\[
 I_{q,k}=\left[e^{i\Phi_q-2\pi iku}
 \left\{\frac{a_q}{ig_k}+\frac{a_q'}{g_k^2}
                   -\frac{a_qg_k'}{g_k^3}\right\}\right]_{a_0}^{b_0}
 +\int_{a_0}^{b_0}\mathcal D_k^2a_q\,
                   e^{i\Phi_q-2\pi iku}\,du.
\]
At a half integer $e^{-2\pi ikz}=(-1)^k$ and $g_k'=-T/u^2$.
The summed boundaries are therefore exactly
\eqref{tr:sublattice-endpoint}.  For $u\ge1/2$, $|k|>M\ge4P$,
\[
 |g_k|\ge\pi|k|,\qquad |g_k'|\le\pi|k|/u,
 \qquad |g_k''|\le2\pi|k|/u^2.
\]
Expansion of the two quotient derivatives gives
\[
 \|\mathcal D_k^2a\|
 \le\frac1{\pi^2k^2}
       \{\|a''\|+3\|a'\|/u+5\|a\|/u^2\}.
\]
The auxiliary range satisfies $q/2\le qu\le3N/2$.
Write $\zeta=-\sigma+(t/2)\log(qu)$.  The reserves imply both
$|\zeta|\le C$ and $-\zeta\ge3/8$.  Direct differentiation gives
\[
 a_q'=\frac{w(qu)}u(\zeta v+v'),\qquad
 a_q''=\frac{w(qu)}{u^2}
 \{(\zeta^2-\zeta+t/2)v+(2\zeta-1)v'+v''\},
\]
and hence
$\|a_q^{(j)}\|\le C_DA_vw(qu)u^{-j}(1+|\rho_{qu}|)^D$ for
$j=0,1,2$.  For $u\ge1$ the weight ratio is at most $u^{-3/8}$;
for $1/2\le u\le1$ it is bounded.  Also
$1+|\rho_{qu}|\le(1+|\rho_q|)(1+|\log u|)$.  Thus
\[
 \int_{1/2}^{N/q+1/2}w(qu)u^{-2}(1+|\rho_{qu}|)^D\,du
 \le C_Dw(q)(1+|\rho_q|)^D.
\]
Summing $\sum_{|k|>M}k^{-2}\le2/M$ proves the claim.  The large
frequency has been integrated into the actual phase, so no factor $T^2$
occurs in the remainder.
\end{proof}

\begin{theorem}[Global coupled primitive transformation]
\label{tr:global-primitive}
Apply \eqref{tr:sublattice-main} to the literal core and block
intervals and set
$\widetilde V_q=\widetilde V_q^{\mathcal C}+\widetilde V_q^{\mathcal B}$.
Define
\[
 \widetilde{\mathcal F}_J=\sum_{q\le N}c_J(q)
 (\operatorname{Re}\widetilde V_q)^TQ
 (\operatorname{Re}\widetilde V_q).
\]
Uniformly for every $1\le J\le N$,
\begin{equation}
\label{tr:primitive-global-bound}
 |\mathcal F_J-\widetilde{\mathcal F}_J|
 \le C_D\|Q\|A_v^2
 \left\{\frac{(Nw_N)L^D}{M}+\frac{L^{2D}}{M^2}\right\}.
\end{equation}
For $M=\lceil4P\rceil\asymp N^2$, this is
\[
 O_D\!\left(\|Q\|A_v^2
       \{L^DN^{-3/2-\kappa/8}+L^{2D}N^{-4}\}\right).
\]
For fixed $D$ and coefficient scales growing by any fixed power of $L$,
it is $o(N^{-\kappa/4})$.
\end{theorem}
\begin{proof}
The reserves imply for $1\le q\le N$
\[
 w(q)\le q^{-s_0},\qquad s_0=35/64>1/2,
\]
since $\sigma-(t/4)\log q\ge c_\sigma+(t/4)\log N\ge35/64$.
The elementary divisor convolution therefore gives
\[
 \sum_{q\le N}\frac{\tau(q)w(q)}q\le\zeta(1+s_0)^2,
 \qquad
 \sum_{q\le N}\tau(q)w(q)^2\le\zeta(2s_0)^2.
\]
Indeed $\sum_{q\ge1}\tau(q)q^{-s}=\sum_{a,b\ge1}(ab)^{-s}
=\zeta(s)^2$ for $s>1$.  No M\"obius cancellation is used.
Let $E_q=\|V_q-\widetilde V_q\|$.  The preceding theorem gives
$E_q\le C_DA_vM^{-1}w(q)(1+|\rho_q|)^D$, while
Equation~\eqref{tr:weighted-moments} gives
$\|V_q\|\le C_DA_v(N/q)w_N$.  The quadratic perturbation inequality is
\[
 \left|(\operatorname{Re}V_q)^TQ\operatorname{Re}V_q
 -(\operatorname{Re}\widetilde V_q)^TQ\operatorname{Re}\widetilde V_q
 \right|\le\|Q\|E_q(2\|V_q\|+E_q).
\]
Multiply by $|c_J(q)|\le\tau(q)$ and sum.  Since
$1+|\rho_q|\le CL$, the two divisor bounds give exactly
\eqref{tr:primitive-global-bound}.  Shared half-integer endpoint terms
cancel when block and core vectors are added; keeping them separately
retains the mixed terms without altering this payment.  The final powers
follow from \eqref{tr:physical-scales}.
\end{proof}

For the original degree-four features the displayed logarithmic powers
are $L^4,L^8$; for the degree-six extension they are $L^6,L^{12}$.
Every matrix and feature scale remains explicit.  As with the product
theorem, a remainder smaller than an available physical upper-budget
scale does not remove the need for the measured physical payment.

\subsection{Resonant sublattices and the remaining signed estimate}

For $k>0$, the phase of $I_{q,k}$ has its unique stationary point at
\[
 u_*=P/k,\qquad n_*=qP/k,
\]
which lies in the open auxiliary interval precisely when
$P/b_0<k<P/a_0$.  Equalities are endpoint transitions.  At the stationary
point the phase is exactly
\[
 \theta_t+T\{\log(qP)-1\}-T\log k.
\]
An isolated interior stationary point with a suitable smooth local
cutoff has leading Gaussian coefficient
\[
 \frac{\sqrt P}{k}w(qP/k)
 v(\log(qP/k)-\mu)
 e^{i\{\theta_t+T(\log(qP)-1)-T\log k-\pi/4\}}.
\]
No replacement of the exact retained integrals by this coefficient was
used in Theorem~\ref{tr:global-primitive}.  Its weight and center have the
exact algebraic reflection
\begin{align*}
 \frac{\sqrt P}{k}w(qP/k)&=C_qk^{-\sigma_q^*}e^{t\log^2k/4},&
 C_q&=\sqrt P(qP)^{-\sigma}e^{t\log^2(qP)/4},\\
 \sigma_q^*&=1-\sigma+\tfrac t2\log(qP),&
 \mu_q^*&=\log(qP)-\mu.
\end{align*}
Thus $\log(qP/k)-\mu=-(\log k-\mu_q^*)$.  Relative to $q=1$,
\[
 \sigma_q^*=\sigma_1^*+\tfrac t2\log q,\qquad
 \mu_q^*=\mu_1^*+\log q,\qquad
 \theta_q^*=\theta_1^*+T\log q,
\]
where $\theta_q^*=\theta_t+T(\log(qP)-1)-\pi/4$.
For fixed $q$ in a macroscopic physical block the modes range
asymptotically from $qN$ to $2qN$, since $P/N^2\to1$; their exact
endpoint ranges remain the half-integer ranges above.  The cutoff is
therefore changed by reflection.  In the ratio term of
\eqref{tr:mobius-formula} the two carriers at a fixed $q$ cancel; in its
product term they remain, including $2T\log q$.  The sublattices cannot
be replaced by independent phases or by a shared $q=1$ reflected sum.

For a disjoint recombination of the complete signed dual, use
\[
 \mathcal K=\mathcal F_J+\mathcal D_{>J}+\mathcal P_{>J},
\]
where the latter terms are exactly the ratio and product contributions
with actual common factor greater than $J$.  The coupled primitive
expression already contains both channels for the small-common-factor
piece.  Adding a separately transformed full product channel to it
would count that product piece twice.  The product theorem may instead
be restricted to the matching $j>J$ intervals.  An arbitrary quadratic
on degree-six features can have total pair degree twelve; its coupled
primitive theorem is valid as stated, but any separate product or
common-factor smoothing invocation must supply the corresponding
fixed-degree amplitude bounds rather than silently invoke the
special degree-six case.

The retained main expression has no proved one-sided sign.  Its
M\"obius coefficients may be negative, and the genuine candidate
conditions constrain only the complete moments.  The remaining
analytical target is a candidate-conditioned upper bound on the coupled
resonant and endpoint terms after this disjoint recombination, strong
enough to beat the transformation, candidate, and physical payments.
The results in these sections neither exclude collisions nor establish
higher-multiplicity or complementary-parameter coverage.

\section{The prescribed leading stationary reflection}
\label{sr:section}
In this section write $a_r=\alpha_r$ and $a_i=\alpha_i$ for the physical
alpha coordinates in \eqref{fo:real-data}.
The following identities concern the leading stationary coefficient,
not a replacement of an endpoint-crossing integral by a full Gaussian.
For the original-index phase $T\log v-2\pi kv$ the positive-frequency
saddle is $v_k=P/k$, where $P=T/(2\pi)$.
It has negative second derivative $-T/v_k^2$, hence leading amplitude
$w(P/k)\sqrt P/k$ and stationary phase
$T\log(P/k)-T-\pi/4$ \cite{dlmfstationary}.

\begin{proposition}[Exact heat-weight reflection]
\label{sr:weight}
For positive $P,k$,
\begin{align}
\frac{\sqrt P}{k}w(P/k)
 &= C_{\rm sd}\,k^{-\sigma^*}e^{t\log^2k/4},\notag\\
\sigma^*&=1-\sigma+\tfrac t2\log P,\qquad
C_{\rm sd}=P^{1/2-\sigma}e^{t\log^2P/4}.
\end{align}
With $\delta_\sigma=\sigma-(1/2+(t/4)\log P)$ this is equivalently
\begin{equation}
\frac{\sqrt P}{k}w(P/k)
 =w(k)e^{\delta_\sigma(2\log k-\log P)} .
\label{sr:amplitude}
\end{equation}
For the physical data,
$\delta_\sigma=O(tx^{-2}+t^2x^{-1})$.
\end{proposition}
\begin{proof}
Expand $\log^2(P/k)=(\log P-\log k)^2$.
Since $\log P=L+\log(1-ta_i/x)$ and $a_r-L/2=O(x^{-2})$,
\[
\delta_\sigma=\frac t2
\{a_r-L/2-\tfrac12\log(1-ta_i/x)\},
\]
which gives the bound. On $c_1\sqrt P\le k\le c_2\sqrt P$,
$2\log k-\log P$ is bounded. The prefactor and power in the
first formula must be combined before bounding them.
\end{proof}

\begin{proposition}[Actual center and carrier cancellation]
\label{sr:reflection}
Set
\[
\Delta_\mu=\log P-2\mu,\qquad
\Theta=2\theta_t+T\log P-T-\pi/4 .
\]
Then $\rho(P/k)=-\rho(k)+\Delta_\mu$, and uniformly for the
physical sector,
\[
\Delta_\mu=O(x^{-2}),\qquad \Theta+2\pi=O(x^{-1}).
\]
\end{proposition}
\begin{proof}
Let $z=ta_i/x$. Using the actual $\mu$,
\[
\Delta_\mu=\log(1-z)-2(a_r-L/2)
                     +z\frac{xV}{1+tU/2}.
\]
Here $xV=1+O(x^{-2})$, $U=O(x^{-2})$, and
$\log(1-z)+z=O(z^2)$. The terms of nominal order $t/x$
therefore cancel.
For the carrier let $T_0=x/2$, $\delta=T-T_0=-ta_i/2$, and
$F(u)=u\log(u/(2\pi))-u$. Taylor's formula gives
\[
F(T)-F(T_0)=\delta L+
\delta^2\int_0^1\frac{1-v}{T_0+v\delta}\dd v.
\]
Substitution of the physical carrier yields exactly
\[
\begin{split}
\Theta+2\pi={}&\tfrac12\arctan x-\tfrac\pi4
 -\frac x4\log(1+x^{-2})+ta_i(a_r-L/2)\\
 &+\delta^2\int_0^1\frac{1-v}{T_0+v\delta}\dd v .
\end{split}
\]
Each term is $O(x^{-1})$ for $t\le1/20$. The nominal $tL$
terms cancel between the thermal carrier and the frequency correction.
\end{proof}

Let $q(v)=w(v)e^{i(\theta_t+T\log v)}$.
The leading carried stationary datum of fixed logarithmic degree $j$ is
exactly
\[
B_j(k)= e^{i\Theta}e^{\delta_\sigma(2\log k-\log P)}
                    \overline{q(k)}\,\{-\rho(k)+\Delta_\mu\}^{j}.
\]
On any fixed macroscopic dual interval
$c_1\sqrt P\le k\le c_2\sqrt P$,
\[
|B_j(k)-(-1)^j\overline{q(k)}\rho(k)^j|\le C_jw(k)/x.
\]
Indeed $\rho(k)$ is bounded there, and the preceding propositions
control the exponential and polynomial differences. There are $O(N)$
indices of weight $O(w_N)$, since $\sqrt P/N\to1$.
The moment comparison costs $O_j(Nw_N/x)$; any fixed quadratic
with coefficient scale $\mathfrak M$ costs
$O(\mathfrak M(Nw_N)^2/x)=O(\mathfrak M N^{-1-\kappa/4})$.
This is a static coefficient comparison on that range, not a global
stationary-phase remainder theorem.

The algebraic reflection $M_j\mapsto(-1)^j\overline{M_j}$
fixes even cosine and odd sine moments. It fixes
$\mathcal K=2Y_3^2+3X_2X_4-\Gamma X_2^2$.
It changes $Z_1=Y_1+\epsilon X_1$ to $Z_1-2\epsilon X_1$,
and, holding the physical drift fixed,
\[
\mathcal J^\dagger(\mathcal R M;\epsilon)
=\mathcal J^\dagger(M;\epsilon)-4\epsilon X_0Y_6
                                      +4\epsilon X_1Y_5
=\mathcal J^\dagger(M;-\epsilon).
\]
The difference has absolute envelope
$O(|\epsilon|(Nw_N)^2)=O(tN^{-1-\kappa/4})$.
An arbitrary reversal of the physical drift is not allowed.

Moreover the saddle maps $[A,B]$ to $[P/B,P/A]$:
$[N/2,N]$ maps approximately to $[N,2N]$, rather than to the
original natural cutoff. Sublattices carry the additional $q$-shifts
already described in the M\"obius section. Thus the useful coefficient
reflection does not identify the complete transformed cutoff with the
original sum and does not turn the threshold quadratic into a negative
square.

\section{The remaining signed implication}
\label{sec:remaining}
One valid exact common-factor split of the extended quadratic is
\[
\mathcal J^\dagger
=\mathcal F_J^\dagger+\mathcal D_{>J}^\dagger
                         +\mathcal P_{>J}^\dagger .
\]
The frontier is coupled, the large-common-factor difference part uses
Theorem~\ref{em:remainder}, and its matching product part uses the finite
Poisson/Fresnel theorem. Alternatively one may transform the entire product
channel and retain the matching difference representation. The two
complete product representations are alternatives; they cannot be added
as disjoint contributions.

If a macroscopic near-pair band is paid separately, subtract that band's
exact contribution from the retained main expression first.
Let $\mathcal J^\dagger_{\rm rem}\le U_{\rm rem}$ be a proved signed
bound for the complement, let $E_{\rm tr}$ be the sum of the errors
of the particular representation used, and let $\mathsf B_c,\mathsf B_s$
be its near-band envelopes. A sufficient ordinary-threshold exclusion is
\begin{equation}
U_{\rm rem}+E_{\rm tr}+\mathsf B_c+\mathsf B_s+\Pi^\dagger
                  +\frac{\widehat\Delta_{\rm count}}{4c^6}<0.
\label{eq:open-target}
\end{equation}
The quantities in this formula include the measured candidate payment
and the physical raw-jet payment with the selected normalizer. Statements
that transformation errors lie below a displayed asymptotic upper-budget
scale do not assert a positive lower bound for a measured error.

The transformations now give explicit analytical main and boundary terms
with vanishing remainders. What is still missing is a one-sided bound for
their complete coupled resonant contribution under the genuine full-sum
candidate constraints. A sign on an isolated block or an independent
sublattice candidate does not imply \eqref{eq:open-target}. Absolute
moment envelopes need not yield its sign. The invariant leading reflection
does not yield it either.

The proposed next analytical step is therefore to seek an inequality for
the coupled stationary integrals and boundary functions with their actual
M\"obius coefficients, shifted centers, thermal weights, and common
frequency. Higher multiplicities and small-curvature branches require
their own tests. Coverage of complementary parameters and the
small-time endpoint is also necessary for a uniform collision theorem.
No assertion in this paper closes those obligations or proves RH.

\section*{Acknowledgements and preparation record}
This manuscript was drafted for Edward Baker with substantial assistance
from OpenAI's Codex large language model. The recorded configuration for
this drafting turn is GPT-6.1-sol, reasoning effort ultra. LLM assistance
included synthesis of research notes, algebraic derivation, exposition,
and internal cross-audits. Those checks are not independent mathematical
validation. Earlier notes retain their own preparation metadata.
The classical heat-flow normalization and imported effective inputs are
credited to the sources below; the analytical reductions are written with
their proof obligations and limits explicit.

\begin{thebibliography}{99}
\bibitem{polymath}
D.\ H.\ J.\ Polymath,
\emph{Effective approximation of heat flow evolution of the Riemann
$\xi$ function, and a new upper bound for the de Bruijn--Newman constant},
arXiv:1904.12438v2 (2019).
\url{https://arxiv.org/abs/1904.12438}.

\bibitem{dlmffourier}
NIST Digital Library of Mathematical Functions, \S1.8,
\emph{Fourier series and Poisson's summation formula}.
\url{https://dlmf.nist.gov/1.8}.

\bibitem{dlmfstationary}
NIST Digital Library of Mathematical Functions, \S2.3(iv),
\emph{Method of stationary phase}.
\url{https://dlmf.nist.gov/2.3#iv}.

\bibitem{dlmfcot}
NIST Digital Library of Mathematical Functions, equation 4.22.3,
\emph{Cotangent partial fractions}.
\url{https://dlmf.nist.gov/4.22.E3}.

\bibitem{dlmfem}
NIST Digital Library of Mathematical Functions, equation 2.10.1,
\emph{Euler--Maclaurin formula}.
\url{https://dlmf.nist.gov/2.10.E1}.

\bibitem{dlmfbern}
NIST Digital Library of Mathematical Functions, equation 24.8.1,
\emph{Bernoulli polynomial Fourier series}.
\url{https://dlmf.nist.gov/24.8.E1}.

\bibitem{projectnotes}
Edward Baker research project, with LLM assistance,
\emph{Prime-phase torus investigation}, internal Notes 1--9,
10 October 2026. Source directory:
\texttt{09\_prime\_phase\_torus/notes/}.
Individual notes state their model, preparation record, and audit scope.

\bibitem{heatnotes}
Edward Baker research project, with LLM assistance,
\emph{Signed shrinking-sector approximation, threshold jets, analytical
Schur cones, density-strengthened thresholds, and stationary reflection},
internal Heat Notes 8, 12, and 20--22, 9--10 October 2026.
Source directory: \texttt{newman\_collisions/notes/}.
\end{thebibliography}
\end{document}
